\documentclass[11pt]{article}

% ============================================================
% LATEX FOR ECON PHD STUDENTS - BEGINNER HANDBOOK
% ============================================================
% Goal: assume zero LaTeX knowledge, but teach the commands and
% habits most useful for economics coursework and paper writing.
% Most examples show BOTH the source code and the compiled result.
% ============================================================

% ---------- Page and typography ----------
\usepackage[margin=0.88in]{geometry}
\usepackage{microtype}
\usepackage{parskip}
\usepackage{tocloft}
\setlength{\cftsecnumwidth}{3.6em}
\setlength{\cftsubsecnumwidth}{4.6em}
\setlength{\parindent}{0pt}

% ---------- Mathematics ----------
\usepackage{amsmath,amssymb,bm,mathtools,amsthm}

% ---------- Tables ----------
\usepackage{booktabs,threeparttable,array,tabularx}
\usepackage{siunitx}

% ---------- Figures ----------
\usepackage{graphicx}
\usepackage{subcaption}
\usepackage{pdflscape}

% ---------- Lists ----------
\usepackage{enumitem}
\setlist{nosep,leftmargin=*}
\setlist[description]{style=nextline,leftmargin=2em,labelindent=0pt}

% ---------- Code examples ----------
\usepackage{xcolor}
\usepackage{listings}
\definecolor{codebg}{RGB}{247,247,247}
\definecolor{rulegray}{RGB}{215,215,215}
\lstset{
  basicstyle=\ttfamily\small,
  backgroundcolor=\color{codebg},
  frame=single,
  rulecolor=\color{rulegray},
  breaklines=true,
  columns=fullflexible,
  showstringspaces=false,
  aboveskip=0.7em,
  belowskip=0.7em,
  xleftmargin=0.2em,
  xrightmargin=0.2em
}

% ---------- Citations and links ----------
\usepackage[round,authoryear]{natbib}
\usepackage[hidelinks]{hyperref}
\usepackage[nameinlink,noabbrev]{cleveref}
\urlstyle{same}

% ---------- Useful layout ----------
\usepackage{float}
\usepackage{placeins}
\usepackage{setspace}
\usepackage{lineno}

% ---------- Custom commands used in examples ----------
\newcommand{\E}{\mathbb{E}}
\newcommand{\Prob}{\mathbb{P}}
\newcommand{\Var}{\operatorname{Var}}
\newcommand{\Cov}{\operatorname{Cov}}
\newcommand{\R}{\mathbb{R}}
\newcommand{\dd}{\,\mathrm{d}}
\newcommand{\code}[1]{\texttt{#1}}

% ---------- Common academic environments ----------
\newtheorem{assumption}{Assumption}
\newtheorem{proposition}{Proposition}
\newtheorem{lemma}{Lemma}
\theoremstyle{definition}
\newtheorem{definition}{Definition}
\newcommand{\sym}[1]{\ifmmode^{#1}\else\(^{#1}\)\fi}

% ---------- Teaching labels ----------
\newcommand{\sourcehead}{\par\smallskip\textbf{What you type}\par}
\newcommand{\resulthead}{\par\smallskip\textbf{What appears in the PDF}\par\smallskip}
\newcommand{\whyhead}{\par\smallskip\textbf{What it means}\par}
\newcommand{\tip}[1]{\par\smallskip\textbf{Beginner tip.} #1\par\smallskip}
\newcommand{\warning}[1]{\par\smallskip\textbf{Common mistake.} #1\par\smallskip}

\title{\textbf{LaTeX for Econ PhD Students}\\[0.35em]
\large A Beginner Handbook for Coursework and Academic Paper Writing}
\author{}
\date{}

\begin{document}
\maketitle

\begin{center}
\begin{minipage}{0.91\textwidth}
\small
\textbf{Who this is for.} You may know economics, econometrics, and mathematics, but know almost nothing about LaTeX. This handbook starts from zero and focuses on the commands and habits you are most likely to use in problem sets, field-course notes, research drafts, and economics papers. It also includes paper-writing patterns that recur in empirical and theoretical economics without turning the handbook into a catalog of econometric formulas.

\textbf{How to use it.} Read Parts I--III once. After that, use the later sections as a lookup reference. You do \emph{not} need to memorize the commands. Copy a working example, change one thing, compile, and repeat.
\end{minipage}
\end{center}

\tableofcontents
\clearpage

% ============================================================
\part{Start Here: How LaTeX Works}
% ============================================================

\section{The One-Minute Mental Model}

LaTeX is a writing system in which you type plain text plus instructions. You save those instructions in a file ending in \code{.tex}. A LaTeX engine reads the file and creates a PDF.

\sourcehead
\begin{lstlisting}
The parameter of interest is $\beta$.
\end{lstlisting}

\resulthead
The parameter of interest is $\beta$.

The command-like parts are not supposed to appear literally in the PDF. They tell LaTeX what to typeset.

A useful mental model is:

\begin{center}
\code{main.tex} $\longrightarrow$ \textbf{compile} $\longrightarrow$ \code{main.pdf}
\end{center}

\tip{Your source file is not the final product. The PDF is the final typeset result.}

\section{Your First File: The Smallest Working Document}

\sourcehead
\begin{lstlisting}
\documentclass[11pt]{article}

\begin{document}

Hello, LaTeX.

\end{document}
\end{lstlisting}

\resulthead
Hello, LaTeX.

\whyhead
\begin{itemize}
  \item \code{\textbackslash documentclass[11pt]\{article\}} chooses a general article layout and an 11-point base font.
  \item \code{\textbackslash begin\{document\}} marks the start of material that will appear in the PDF.
  \item \code{\textbackslash end\{document\}} marks the end.
\end{itemize}

Everything before \code{\textbackslash begin\{document\}} is called the \textbf{preamble}. Packages, custom commands, title information, and document settings normally go there.

\section{What Does ``Compile'' Mean?}

Compiling means asking LaTeX to turn the \code{.tex} source into a PDF. In Overleaf you press \textbf{Recompile}. In a local editor such as VS Code, TeXstudio, or RStudio, you run a LaTeX build command.

A simple project might contain:

\begin{lstlisting}
main.tex          your paper source
main.pdf          compiled paper
main.aux          temporary cross-reference information
main.log          compiler log
main.out          hyperlink information
references.bib    bibliography database (if used)
\end{lstlisting}

You normally edit \code{main.tex}. The auxiliary files are generated automatically.

\subsection{Why do I sometimes need to compile twice?}

LaTeX may need one pass to discover numbers and a second pass to insert them. This is common for:

\begin{itemize}
  \item section, equation, figure, and table references;
  \item table of contents entries;
  \item citations and bibliographies.
\end{itemize}

If a reference appears as \textbf{??}, compile again before assuming something is wrong.

\section{Commands, Arguments, Options, and Environments}

Most LaTeX syntax falls into four patterns.

\subsection{A command}

\sourcehead
\begin{lstlisting}
\LaTeX
\end{lstlisting}
\resulthead
\LaTeX

A command begins with a backslash.

\subsection{A command with a required argument}

\sourcehead
\begin{lstlisting}
\textbf{Important result}
\end{lstlisting}
\resulthead
\textbf{Important result}

The braces \code{\{...\}} contain the material the command acts on.

\subsection{A command with an optional argument}

\sourcehead
\begin{lstlisting}
\documentclass[11pt]{article}
\end{lstlisting}

Square brackets usually contain optional settings; braces usually contain required information.

\subsection{An environment}

\sourcehead
\begin{lstlisting}
\begin{itemize}
  \item First point
  \item Second point
\end{itemize}
\end{lstlisting}

\resulthead
\begin{itemize}
  \item First point
  \item Second point
\end{itemize}

Every \code{\textbackslash begin\{name\}} must eventually have a matching \code{\textbackslash end\{name\}}.

\section{Packages: Adding Features}

Packages add functionality. A practical beginner preamble for economics paper writing is:

\begin{lstlisting}
\documentclass[11pt]{article}

\usepackage[margin=1in]{geometry}
\usepackage{amsmath,amssymb}
\usepackage{graphicx}
\usepackage{subcaption}
\usepackage{pdflscape}
\usepackage{booktabs}
\usepackage{threeparttable}
\usepackage[round,authoryear]{natbib}
\usepackage[hidelinks]{hyperref}
\end{lstlisting}

\begin{center}
\begin{tabularx}{0.94\textwidth}{@{}lX@{}}
\toprule
Package & Why you may want it \\
\midrule
\code{geometry} & Set margins simply. \\
\code{amsmath} & Better equations and environments such as \code{align}. \\
\code{amssymb} & Extra mathematical symbols. \\
\code{graphicx} & Insert PDF/PNG/JPG figures. \\
\code{booktabs} & Professional horizontal rules for tables. \\
\code{threeparttable} & Clean notes below tables. \\
\code{natbib} & Author-year citations such as ``Author (2020)'' and ``(Author, 2020).'' \\
\code{hyperref} & Clickable cross-references and links. \\
\bottomrule
\end{tabularx}
\end{center}

\tip{Start with a small preamble. Add packages when you actually need them. A huge preamble makes debugging harder.}

% ============================================================
\part{Text, Punctuation, and Document Structure}
% ============================================================

\section{Plain Text, Paragraphs, and Line Breaks}

You type ordinary prose normally.

\sourcehead
\begin{lstlisting}
This paper studies how households respond to a change in prices.
\end{lstlisting}
\resulthead
This paper studies how households respond to a change in prices.

To create a new paragraph, leave a blank line in the source.

\sourcehead
\begin{lstlisting}
This is the first paragraph.

This is the second paragraph.
\end{lstlisting}

\resulthead
This is the first paragraph.

This is the second paragraph.

\warning{Do not use repeated \code{\textbackslash\textbackslash} commands to create ordinary paragraph spacing. A blank source line is the normal paragraph break.}

\section{Bold, Italic, Emphasis, and Typewriter Text}

\begin{center}
\begin{tabularx}{0.94\textwidth}{@{}lXX@{}}
\toprule
Purpose & Source & Compiled result \\
\midrule
Bold & \code{\textbackslash textbf\{important\}} & \textbf{important} \\
Italic & \code{\textbackslash textit\{term\}} & \textit{term} \\
Emphasis & \code{\textbackslash emph\{key idea\}} & \emph{key idea} \\
Code/file name & \code{\textbackslash texttt\{main.tex\}} & \texttt{main.tex} \\
\bottomrule
\end{tabularx}
\end{center}

\tip{In academic papers, use bold and italics sparingly. Section structure should do most of the organizational work.}

\section{Reserved Characters: Characters LaTeX Treats Specially}

Some keyboard characters are instructions to LaTeX. If you want the literal character in prose, you must escape or replace it.

\begin{center}
\begin{tabularx}{0.96\textwidth}{@{}c l X@{}}
\toprule
Character & Type this to print it & Why it is special \\
\midrule
\% & \code{\textbackslash\%} & Starts a comment. \\
\$ & \code{\textbackslash\$} & Starts/ends inline math mode. \\
\& & \code{\textbackslash\&} & Separates columns/alignment points. \\
\_ & \code{\textbackslash\_} & Creates a subscript in math mode. \\
\# & \code{\textbackslash\#} & Used for macro parameters. \\
\{ & \code{\textbackslash\{} & Starts a grouped argument. \\
\} & \code{\textbackslash\}} & Ends a grouped argument. \\
\textasciitilde & \code{\textbackslash textasciitilde} & \code{\~} is an accent command; bare \code{~} is a nonbreaking space. \\
\textasciicircum & \code{\textbackslash textasciicircum} & Superscript operator in math mode. \\
\textbackslash & \code{\textbackslash textbackslash} & Starts most LaTeX commands. \\
\bottomrule
\end{tabularx}
\end{center}

\sourcehead
\begin{lstlisting}
R\&D spending increased by 5\%.
The replication file is data\_clean.csv.
The fee is \$25.
See Figure \#2.
\end{lstlisting}

\resulthead
R\&D spending increased by 5\%.  The replication file is data\_clean.csv. The fee is \$25. See Figure \#2.

\warning{Typing \code{5\%} without the backslash makes LaTeX treat the rest of that source line as a comment.}

\section{Comments: Notes That Do Not Print}

Anything after an unescaped percent sign is ignored until the end of that source line.

\sourcehead
\begin{lstlisting}
% TODO: rewrite this paragraph.
This sentence appears in the PDF. % private note to myself
\end{lstlisting}

\resulthead
This sentence appears in the PDF.

Comments are useful for TODOs, temporarily disabling text, and organizing long source files.

\section{Quotation Marks, Dashes, and Small Typography Habits}

Traditional LaTeX quotation marks are typed with two backticks and two apostrophes.

\sourcehead
\begin{lstlisting}
``Identification requires a stronger assumption.''
\end{lstlisting}
\resulthead
``Identification requires a stronger assumption.''

Use different numbers of hyphens for different dashes:

\begin{center}
\begin{tabular}{lll}
\toprule
Source & Result & Typical use \\
\midrule
\code{1-5} & 1-5 & hyphen / simple connection \\
\code{1990--2000} & 1990--2000 & numerical range (en dash) \\
\code{word---word} & word---word & parenthetical break (em dash) \\
\bottomrule
\end{tabular}
\end{center}

For references, a tilde creates a nonbreaking space:

\sourcehead
\begin{lstlisting}
Table~\ref{tab:main}
Figure~\ref{fig:main}
Section~\ref{sec:data}
\end{lstlisting}

This keeps ``Table'' and its number on the same line.

\section{Titles, Authors, Dates, and Abstracts}

\sourcehead
\begin{lstlisting}
\title{The Effect of X on Y}
\author{Your Name\\Your University}
\date{September 2026}

\begin{document}
\maketitle

\begin{abstract}
This paper studies ...
\end{abstract}
\end{lstlisting}

\resulthead
\begin{center}
{\Large\bfseries The Effect of X on Y}\par
\vspace{0.5em}
Your Name\par
Your University\par
\vspace{0.4em}
September 2026
\end{center}

\begin{abstract}
This paper studies ...
\end{abstract}

If you do not want a date, use \code{\textbackslash date\{\}}.

\section{Sections and Subsections}

\sourcehead
\begin{lstlisting}
\section{Introduction}
\subsection{Data Sources}
\subsubsection{Sample Construction}
\end{lstlisting}

LaTeX handles the numbering automatically. Do not type ``1. Introduction'' yourself.

For an unnumbered heading:

\begin{lstlisting}
\section*{Acknowledgments}
\end{lstlisting}

\section{A Beginner-Friendly Economics Paper Structure}

A common empirical-paper skeleton is:

\begin{lstlisting}
\section{Introduction}
\section{Background or Institutional Setting}
\section{Data}
\section{Research Design or Empirical Strategy}
\section{Results}
\section{Robustness and Extensions}
\section{Conclusion}
\end{lstlisting}

The exact headings depend on the paper. LaTeX does not determine your research structure; it only typesets it consistently.

% ============================================================
\part{Math from Zero: Syntax You Will Reuse}
% ============================================================

\section{Math Mode: The Most Important Distinction}

LaTeX treats text and mathematics differently.

\subsection{Inline math}

\sourcehead
\begin{lstlisting}
The coefficient of interest is $\beta$.
\end{lstlisting}
\resulthead
The coefficient of interest is $\beta$.

Use inline math for short expressions inside a sentence.

\subsection{Displayed math}

\sourcehead
\begin{lstlisting}
\[
y_i = \alpha + \beta x_i + \varepsilon_i.
\]
\end{lstlisting}
\resulthead
\[
y_i = \alpha + \beta x_i + \varepsilon_i.
\]

\warning{Do not write math commands such as \code{\textbackslash beta} in ordinary text mode. Write \code{\$\textbackslash beta\$}.}

\section{Subscripts and Superscripts}

\sourcehead
\begin{lstlisting}
$x_i$
$x_{it}$
$x^2$
$x^{10}$
$x_{it}^2$
\end{lstlisting}

\resulthead
$x_i$, \quad $x_{it}$, \quad $x^2$, \quad $x^{10}$, \quad $x_{it}^2$.

Use braces whenever the subscript or superscript contains more than one token.

\warning{\code{x\_it} does not mean $x_{it}$. Write \code{x\_\{it\}}.}

\section{Greek Letters}

\begin{center}
\begin{tabularx}{0.94\textwidth}{@{}llll@{}}
\toprule
Source & Result & Source & Result \\
\midrule
\code{\textbackslash alpha} & $\alpha$ & \code{\textbackslash beta} & $\beta$ \\
\code{\textbackslash gamma} & $\gamma$ & \code{\textbackslash delta} & $\delta$ \\
\code{\textbackslash epsilon} & $\epsilon$ & \code{\textbackslash varepsilon} & $\varepsilon$ \\
\code{\textbackslash theta} & $\theta$ & \code{\textbackslash lambda} & $\lambda$ \\
\code{\textbackslash mu} & $\mu$ & \code{\textbackslash pi} & $\pi$ \\
\code{\textbackslash rho} & $\rho$ & \code{\textbackslash sigma} & $\sigma$ \\
\code{\textbackslash phi} & $\phi$ & \code{\textbackslash omega} & $\omega$ \\
\code{\textbackslash Delta} & $\Delta$ & \code{\textbackslash Sigma} & $\Sigma$ \\
\bottomrule
\end{tabularx}
\end{center}

\section{Fractions, Roots, Parentheses, and Brackets}

\sourcehead
\begin{lstlisting}
$\frac{a}{b}$
$\sqrt{x}$
$\sqrt[3]{x}$
\end{lstlisting}
\resulthead
$\frac{a}{b}$, \quad $\sqrt{x}$, \quad $\sqrt[3]{x}$.

For tall expressions, use automatically sized delimiters:

\sourcehead
\begin{lstlisting}
\[
\left( \frac{a+b}{c} \right)
\qquad
\left[ \frac{a+b}{c} \right]
\]
\end{lstlisting}
\resulthead
\[
\left( \frac{a+b}{c} \right)
\qquad
\left[ \frac{a+b}{c} \right].
\]

Literal curly braces inside math need backslashes:

\sourcehead
\begin{lstlisting}
\[
\left\{x \in \mathbb{R}: x>0\right\}
\]
\end{lstlisting}
\resulthead
\[
\left\{x \in \mathbb{R}: x>0\right\}.
\]

\section{Common Mathematical Symbols}

\begin{center}
\begin{tabularx}{0.95\textwidth}{@{}lcl lcl@{}}
\toprule
Source & Result & Meaning & Source & Result & Meaning \\
\midrule
\code{\textbackslash le} & $\le$ & less/equal & \code{\textbackslash ge} & $\ge$ & greater/equal \\
\code{\textbackslash neq} & $\neq$ & not equal & \code{\textbackslash approx} & $\approx$ & approximately \\
\code{\textbackslash in} & $\in$ & belongs to & \code{\textbackslash notin} & $\notin$ & not in \\
\code{\textbackslash infty} & $\infty$ & infinity & \code{\textbackslash pm} & $\pm$ & plus/minus \\
\code{\textbackslash to} & $\to$ & tends/goes to & \code{\textbackslash implies} & $\implies$ & implies \\
\code{\textbackslash times} & $\times$ & multiplication & \code{\textbackslash cdot} & $\cdot$ & centered dot \\
\bottomrule
\end{tabularx}
\end{center}

\section{Sums, Products, Integrals, and Limits}

\sourcehead
\begin{lstlisting}
\[
\sum_{i=1}^{n} x_i
\qquad
\prod_{i=1}^{n} x_i
\]

\[
\int_a^b f(x)\,\mathrm{d}x
\qquad
\lim_{n\to\infty} a_n
\]
\end{lstlisting}

\resulthead
\[
\sum_{i=1}^{n}x_i
\qquad
\prod_{i=1}^{n}x_i
\]
\[
\int_a^b f(x)\,\mathrm{d}x
\qquad
\lim_{n\to\infty}a_n.
\]

The pattern is reusable: \code{\_\{lower\}} puts something below, and \code{\textasciicircum\{upper\}} puts something above.

\section{Derivatives and Partial Derivatives}

\sourcehead
\begin{lstlisting}
\[
\frac{\mathrm{d}y}{\mathrm{d}x}
\qquad
\frac{\partial f}{\partial x}
\qquad
\frac{\partial^2 f}{\partial x^2}
\]
\end{lstlisting}
\resulthead
\[
\frac{\mathrm{d}y}{\mathrm{d}x}
\qquad
\frac{\partial f}{\partial x}
\qquad
\frac{\partial^2 f}{\partial x^2}.
\]

\section{Math Operators: Use Commands Instead of Italic Letters}

Commands such as \code{\textbackslash log}, \code{\textbackslash ln}, \code{\textbackslash exp}, \code{\textbackslash max}, and \code{\textbackslash min} give proper operator spacing.

\sourcehead
\begin{lstlisting}
\[
\log x,\quad \ln x,\quad \exp(x),\quad
\max_x f(x)
\]
\end{lstlisting}
\resulthead
\[
\log x,\quad \ln x,\quad \exp(x),\quad \max_x f(x).
\]

\warning{Writing \code{max} directly in math produces italic letters $max$, not the mathematical operator $\max$.}

\section{Expectations, Probability, Variance, and Covariance}

\sourcehead
\begin{lstlisting}
\[
\mathbb{E}[Y\mid X]
\qquad
\mathbb{P}(A)
\qquad
\operatorname{Var}(X)
\qquad
\operatorname{Cov}(X,Y)
\]
\end{lstlisting}
\resulthead
\[
\mathbb{E}[Y\mid X]
\qquad
\mathbb{P}(A)
\qquad
\operatorname{Var}(X)
\qquad
\operatorname{Cov}(X,Y).
\]

If you use these constantly, define shortcuts in the preamble:

\begin{lstlisting}
\newcommand{\E}{\mathbb{E}}
\newcommand{\Prob}{\mathbb{P}}
\newcommand{\Var}{\operatorname{Var}}
\newcommand{\Cov}{\operatorname{Cov}}
\end{lstlisting}

Then \code{\textbackslash E[Y\textbackslash mid X]} produces $\E[Y\mid X]$.

\section{Hats, Bars, Tildes, Stars, and Bold Symbols}

\begin{center}
\begin{tabularx}{0.94\textwidth}{@{}lclX@{}}
\toprule
Source & Result & Typical reading & Comment \\
\midrule
\code{\textbackslash hat\{\textbackslash beta\}} & $\hat\beta$ & beta hat & estimate or fitted quantity \\
\code{\textbackslash bar\{x\}} & $\bar x$ & x bar & average / transformed notation \\
\code{\textbackslash tilde\{x\}} & $\tilde x$ & x tilde & transformed/alternative variable \\
\code{x\^*} & $x^*$ & x star & optimum/equilibrium/etc. \\
\code{\textbackslash mathbf\{x\}} & $\mathbf{x}$ & bold x & vector/matrix notation \\
\code{\textbackslash bm\{\textbackslash beta\}} & $\bm{\beta}$ & bold beta & bold Greek symbol \\
\bottomrule
\end{tabularx}
\end{center}

\section{Matrices and Vectors}

\sourcehead
\begin{lstlisting}
\[
\mathbf{x}=
\begin{pmatrix}
x_1\\
x_2\\
x_3
\end{pmatrix}
\qquad
A=
\begin{pmatrix}
a_{11} & a_{12}\\
a_{21} & a_{22}
\end{pmatrix}
\]
\end{lstlisting}

\resulthead
\[
\mathbf{x}=
\begin{pmatrix}
x_1\\x_2\\x_3
\end{pmatrix}
\qquad
A=
\begin{pmatrix}
a_{11} & a_{12}\\a_{21} & a_{22}
\end{pmatrix}.
\]

Inside matrices, \code{\&} separates columns and \code{\textbackslash\textbackslash} ends a row.

\section{Cases and Text Inside Math}

\sourcehead
\begin{lstlisting}
\[
D_i=
\begin{cases}
1, & \text{if unit $i$ is treated},\\
0, & \text{otherwise}.
\end{cases}
\]
\end{lstlisting}

\resulthead
\[
D_i=
\begin{cases}
1, & \text{if unit $i$ is treated},\\
0, & \text{otherwise}.
\end{cases}
\]

Use \code{\textbackslash text\{...\}} for ordinary words inside math mode.

% ============================================================
\part{Equations in an Academic Paper}
% ============================================================

\section{Unnumbered vs. Numbered Equations}

Use \code{\textbackslash[ ... \textbackslash]} when you do not need to refer to the equation later.

Use the \code{equation} environment when the equation deserves a number.

\sourcehead
\begin{lstlisting}
\begin{equation}
y_i = \alpha + \beta x_i + \varepsilon_i.
\label{eq:baseline}
\end{equation}
\end{lstlisting}

\resulthead
\begin{equation}
y_i = \alpha + \beta x_i + \varepsilon_i.
\label{eq:baseline-demo}
\end{equation}

Then write:

\sourcehead
\begin{lstlisting}
Equation~\eqref{eq:baseline} defines the baseline specification.
\end{lstlisting}

\resulthead
Equation~\eqref{eq:baseline-demo} defines the baseline specification.

\section{Multi-Line Equations with align}

The \code{align} environment is one of the most useful tools in economics.

\sourcehead
\begin{lstlisting}
\begin{align}
A
&= B+C,\\
&= D+E,\\
&= F.
\end{align}
\end{lstlisting}

\resulthead
\begin{align}
A
&= B+C,\\
&= D+E,\\
&= F.
\end{align}

The \code{\&} marks the alignment point. Usually you place it before the equals sign.

For no equation numbers, use \code{align*}.

\warning{Inside \code{align}, a literal ampersand is not ``and.'' It controls alignment.}

\section{Long Equations: split and aligned}

If one numbered equation is too wide, split it over multiple lines while keeping one equation number.

\sourcehead
\begin{lstlisting}
\begin{equation}
\begin{split}
y_i
&= \alpha + \beta_1 x_{1i} + \beta_2 x_{2i} \\
&\quad + \beta_3 x_{3i} + \varepsilon_i.
\end{split}
\label{eq:long}
\end{equation}
\end{lstlisting}

\resulthead
\begin{equation}
\begin{split}
y_i
&= \alpha + \beta_1 x_{1i} + \beta_2 x_{2i} \\
&\quad + \beta_3 x_{3i} + \varepsilon_i.
\end{split}
\label{eq:long-demo}
\end{equation}

\section{Equation Punctuation: Equations Are Part of Sentences}

A displayed equation usually belongs grammatically to the sentence around it. Punctuate it accordingly.

\sourcehead
\begin{lstlisting}
We define the transformed outcome as
\[
\tilde y_i = y_i-\bar y,
\]
where $\bar y$ denotes the sample mean.
\end{lstlisting}

\resulthead
We define the transformed outcome as
\[
\tilde y_i = y_i-\bar y,
\]
where $\bar y$ denotes the sample mean.

Notice the comma after the equation because the sentence continues.

\section{Explain Notation After an Equation}

Do not display a complicated equation and leave the reader to infer every symbol.

\sourcehead
\begin{lstlisting}
We write the estimating equation as
\begin{equation}
y_{it}=\alpha+\beta d_{it}+\mathbf{x}_{it}'\gamma+\varepsilon_{it}.
\label{eq:model}
\end{equation}
Here, $y_{it}$ is the outcome, $d_{it}$ is the variable of interest,
$\mathbf{x}_{it}$ is a vector of controls, and $\varepsilon_{it}$ is an error term.
\end{lstlisting}

\tip{This is a writing habit, not just a LaTeX habit: define notation when the reader first needs it.}

\section{Useful Spacing Commands in Math}

Most spacing should be automatic, but a few commands are useful:

\begin{center}
\begin{tabular}{lll}
\toprule
Command & Size & Example \\
\midrule
\code{\textbackslash,} & small space & $f(x)\,\mathrm{d}x$ \\
\code{\textbackslash quad} & larger space & $A\quad B$ \\
\code{\textbackslash qquad} & even larger & $A\qquad B$ \\
\code{\textbackslash!} & small negative space & occasionally tightens notation \\
\bottomrule
\end{tabular}
\end{center}

Do not manually space every operator. Let LaTeX do most of the work.

% ============================================================
\part{Small Commands and Layout Tricks You Will Use Constantly}
% ============================================================

\section{Why These ``Tiny'' Commands Matter}

Once you know basic LaTeX, much of the day-to-day improvement in a paper comes from small commands rather than new environments. Commands such as \code{\textbackslash qquad}, \code{\textasciitilde}, \code{\textbackslash hfill}, \code{\textbackslash noindent}, and \code{\textbackslash notag} help you control spacing, alignment, line breaking, references, and mathematical typography without manually fighting the page.

A useful way to remember them is by job:

\begin{center}
\begin{tabularx}{0.96\textwidth}{lX}
\toprule
Job & Commands worth recognizing \\
\midrule
Small math spacing & \code{\textbackslash,}, \code{\textbackslash:}, \code{\textbackslash;}, \code{\textbackslash!}, \code{\textbackslash quad}, \code{\textbackslash qquad} \\
Horizontal layout & \code{\textbackslash hspace}, \code{\textbackslash hfill}, \code{\textbackslash makebox}, \code{\textbackslash phantom} \\
Vertical layout & \code{\textbackslash smallskip}, \code{\textbackslash medskip}, \code{\textbackslash bigskip}, \code{\textbackslash vspace} \\
Paragraphs and alignment & blank line, \code{\textbackslash par}, \code{\textbackslash noindent}, \code{\textbackslash centering}, \code{\textbackslash raggedright} \\
Equation alignment & \code{\&}, \code{\textbackslash\textbackslash}, \code{\textbackslash notag}, \code{\textbackslash intertext}, \code{aligned} \\
Math typography & \code{\textbackslash text}, \code{\textbackslash mathrm}, \code{\textbackslash mathbf}, \code{\textbackslash bm}, \code{\textbackslash mathbb}, \code{\textbackslash operatorname} \\
Delimiter sizing & \code{\textbackslash left}, \code{\textbackslash right}, \code{\textbackslash bigl}, \code{\textbackslash Bigr} \\
References and links & \code{\textasciitilde}, \code{\textbackslash ref}, \code{\textbackslash eqref}, \code{\textbackslash pageref}, \code{\textbackslash autoref}, \code{\textbackslash href} \\
Page control & \code{\textbackslash newpage}, \code{\textbackslash clearpage}, \code{\textbackslash pagebreak}, \code{\textbackslash FloatBarrier} \\
Reuse & \code{\textbackslash input}, \code{\textbackslash newcommand}, \code{\textbackslash ensuremath} \\
\bottomrule
\end{tabularx}
\end{center}

\tip{Do not try to memorize this table. The goal is recognition: when you need a little more space, a nonbreaking reference, an unnumbered line, or a reusable abbreviation, you should know that LaTeX probably has a command for it.}

\section{Horizontal Spacing in Math: From Tiny to Large}

LaTeX already spaces mathematical operators well. Manual spacing is mainly for semantic separation: for example, before a differential, between a displayed equation and an indexing condition, or inside a compact piece of notation.

\sourcehead
\begin{lstlisting}
$a\!b$       % negative thin space
$a\,b$       % thin space
$a\:b$       % medium space
$a\;b$       % thick space
$a\quad b$   % 1 em
$a\qquad b$  % 2 em
\end{lstlisting}

\resulthead
\[
a\!b \qquad
 a\,b \qquad
 a\:b \qquad
 a\;b \qquad
 a\quad b \qquad
 a\qquad b
\]

The practical hierarchy is
\[
\text{negative: }\! \quad < \quad \text{small: }\, \quad < \quad \: \quad < \quad \; \quad < \quad \quad \quad < \quad \qquad.
\]

\subsection{Economics-paper examples}

A differential often looks cleaner with a thin space:

\sourcehead
\begin{lstlisting}
\int_0^1 f(x)\,\mathrm{d}x
\end{lstlisting}
\resulthead
\[
\int_0^1 f(x)\,\mathrm{d}x
\]

A condition can be separated from the main equation with \code{\textbackslash qquad}:

\sourcehead
\begin{lstlisting}
y_i = \alpha + \beta x_i + \varepsilon_i,
\qquad i=1,\ldots,n.
\end{lstlisting}
\resulthead
\[
y_i = \alpha + \beta x_i + \varepsilon_i,
\qquad i=1,\ldots,n.
\]

A negative thin space is sometimes useful when two symbols look too far apart:

\sourcehead
\begin{lstlisting}
\int\!\!\int f(x,y)\,\mathrm{d}x\,\mathrm{d}y
\end{lstlisting}
\resulthead
\[
\int\!\!\int f(x,y)\,\mathrm{d}x\,\mathrm{d}y
\]

\warning{Do not insert \code{\textbackslash quad} everywhere to make an equation ``look nicer.'' Good mathematical spacing should usually communicate structure, not compensate for a poorly organized expression.}

\section{Fixed and Flexible Horizontal Space in Text}

\subsection{hspace: insert a specific amount of space}

\sourcehead
\begin{lstlisting}
Left\hspace{2em}Right
\end{lstlisting}
\resulthead
Left\hspace{2em}Right

Common units include \code{em} (roughly the width of the current letter M), \code{ex} (related to the current x-height), \code{pt}, \code{cm}, and fractions of text width such as \code{0.2\textbackslash textwidth}.

\tip{Inside prose, manual \code{\textbackslash hspace} is rarely the right way to align material. Tables, minipages, or \code{\textbackslash hfill} are usually more robust.}

\subsection{hfill: stretch space as much as possible}

\sourcehead
\begin{lstlisting}
Draft version \hfill September 2026
\end{lstlisting}
\resulthead
Draft version \hfill September 2026

This is useful when you want material on opposite sides of the same line.

\sourcehead
\begin{lstlisting}
Ning Zhang \hfill Washington State University
\end{lstlisting}
\resulthead
Ning Zhang \hfill Washington State University

\subsection{makebox: reserve a box of a chosen width}

\sourcehead
\begin{lstlisting}
\makebox[3cm][l]{Coefficient:} 0.125
\makebox[3cm][l]{Std. error:}  0.041
\end{lstlisting}
\resulthead
\makebox[3cm][l]{Coefficient:} 0.125\\
\makebox[3cm][l]{Std. error:} 0.041

The second optional argument controls alignment: \code{l}, \code{c}, or \code{r}.

\warning{For a real results table, use \code{tabular} or an imported table rather than manually aligning rows with \code{\textbackslash makebox}. This command is mainly useful for small layout tasks.}

\section{Nonbreaking Spaces: The Tiny Tilde You Will Use Everywhere}

In ordinary text, a space permits a line break. The character \code{\textasciitilde} creates a space but forbids a line break at that location.

\sourcehead
\begin{lstlisting}
Table~\ref{tab:main}
Figure~\ref{fig:event}
Section~\ref{sec:data}
Appendix~\ref{app:proofs}
Equation~\eqref{eq:baseline}
\end{lstlisting}

This prevents awkward breaks such as

\begin{quote}
Table\\
3
\end{quote}

and keeps ``Table 3'' together.

\subsection{Other useful nonbreaking combinations}

\sourcehead
\begin{lstlisting}
pp.~12--15
Figure~A1
$10~\mathrm{km}$
\$5~million
\end{lstlisting}

\resulthead
pp.~12--15 \qquad Figure~A1 \qquad $10~\mathrm{km}$ \qquad \$5~million

\tip{A simple habit: whenever a word and the number that follows it should stay together, consider \code{\textasciitilde}. This is especially useful for Table, Figure, Section, Appendix, Equation, Proposition, and page references.}

\section{Vertical Spacing: smallskip, medskip, bigskip, and vspace}

For small manual adjustments, LaTeX provides named vertical skips:

\sourcehead
\begin{lstlisting}
First paragraph.

\smallskip
A little extra space.

\medskip
A medium amount of extra space.

\bigskip
A larger amount of extra space.
\end{lstlisting}

For an exact amount, use \code{\textbackslash vspace}:

\sourcehead
\begin{lstlisting}
\vspace{0.5em}
\vspace{6pt}
\vspace{0.25in}
\end{lstlisting}

The starred form \code{\textbackslash vspace*\{...\}} also works at the top of a page, where ordinary \code{\textbackslash vspace} may disappear.

\warning{If you find yourself using many \code{\textbackslash vspace} commands throughout a paper, the real problem is probably the document style, float placement, or a table/figure design. Manual vertical spacing should be occasional.}

\section{Paragraphs, Indentation, and Local Alignment}

\subsection{Start a new paragraph with a blank line}

\sourcehead
\begin{lstlisting}
This is the first paragraph.

This is the second paragraph.
\end{lstlisting}

A blank line is equivalent to \code{\textbackslash par}:

\sourcehead
\begin{lstlisting}
This is the first paragraph.\par
This is the second paragraph.
\end{lstlisting}

\subsection{noindent and indent}

\sourcehead
\begin{lstlisting}
\noindent This paragraph does not begin with an indent.

\indent This paragraph explicitly begins with an indent.
\end{lstlisting}

Because this handbook uses the \code{parskip} package, its paragraphs are separated by vertical space and normally have no indentation. Many journal styles instead use indentation and little or no paragraph spacing.

\subsection{centering, raggedright, and raggedleft}

Inside an environment such as a figure, \code{\textbackslash centering} is usually preferable to opening another \code{center} environment.

\sourcehead
\begin{lstlisting}
\begin{figure}[htbp]
  \centering
  \includegraphics[width=.7\textwidth]{figures/main.pdf}
  \caption{Main result}
\end{figure}
\end{lstlisting}

For local text alignment, group the declaration in braces:

\sourcehead
\begin{lstlisting}
{\raggedright Left-aligned text without full justification.\par}
{\centering Centered text.\par}
{\raggedleft Right-aligned text.\par}
\end{lstlisting}

\tip{Commands such as \code{\textbackslash centering} are declarations: they remain active until the current group or environment ends. Braces are a convenient way to limit their scope.}

\section{Alignment Inside Equations: Ampersands, Line Breaks, notag, and intertext}

In \code{align}, the ampersand \code{\&} marks the alignment point. Most economics papers align at the equality sign.

\sourcehead
\begin{lstlisting}
\begin{align}
y_i &= \alpha + \beta x_i + \varepsilon_i, \\
\E[\varepsilon_i\mid x_i] &= 0.
\end{align}
\end{lstlisting}

\resulthead
\begin{align}
y_i &= \alpha + \beta x_i + \varepsilon_i, \\
\E[\varepsilon_i\mid x_i] &= 0.
\end{align}

Use \code{\textbackslash notag} when one line should not receive an equation number:

\sourcehead
\begin{lstlisting}
\begin{align}
A &= B + C \notag\\
  &= D + E.
\end{align}
\end{lstlisting}

\resulthead
\begin{align}
A &= B + C \notag\\
  &= D + E.
\end{align}

Use \code{\textbackslash intertext\{...\}} when a sentence belongs between aligned lines:

\sourcehead
\begin{lstlisting}
\begin{align}
A &= B + C,\\
\intertext{Substituting the definition of $C$ gives}
A &= B + D + E.
\end{align}
\end{lstlisting}

\resulthead
\begin{align}
A &= B + C,\\
\intertext{Substituting the definition of $C$ gives}
A &= B + D + E.
\end{align}

\subsection{aligned: alignment inside another display}

\sourcehead
\begin{lstlisting}
\[
\begin{aligned}
A &= B+C,\\
  &= D+E.
\end{aligned}
\]
\end{lstlisting}

\tip{Use \code{align} for a normal multi-line numbered display. Use \code{aligned} when the aligned block must live inside another math environment, such as inside brackets or one larger equation.}

\section{Invisible Space: phantom, hphantom, and vphantom}

A phantom occupies the same space as its argument but prints nothing.

\sourcehead
\begin{lstlisting}
A = 1.25

\hphantom{A = }0.75
\end{lstlisting}

\resulthead
A = 1.25\\
\hphantom{A = }0.75

\code{\textbackslash hphantom} reserves only horizontal space; \code{\textbackslash vphantom} reserves only vertical space; \code{\textbackslash phantom} reserves both.

A useful mathematical example is forcing delimiters to have matching height:

\sourcehead
\begin{lstlisting}
\left( x_i \vphantom{\frac{A}{B}} \right)
\qquad
\left( \frac{A}{B} \right)
\end{lstlisting}

\resulthead
\[
\left( x_i \vphantom{\frac{A}{B}} \right)
\qquad
\left( \frac{A}{B} \right)
\]

\tip{Phantoms are a fine finishing tool, but if you are using many of them to align a table, switch to a proper table-column specification instead.}

\section{Text Inside Math: text, mathrm, operatorname, and Friends}

Mathematical letters are italic by default because LaTeX assumes they are variables. Words and named operators need different commands.

\subsection{text: ordinary words inside an equation}

\sourcehead
\begin{lstlisting}
y_i =
\begin{cases}
1, & \text{if employed},\\
0, & \text{otherwise}.
\end{cases}
\end{lstlisting}

\resulthead
\[
y_i =
\begin{cases}
1, & \text{if employed},\\
0, & \text{otherwise}.
\end{cases}
\]

\subsection{mathrm: upright mathematical labels}

\sourcehead
\begin{lstlisting}
GDP_{it}^{\mathrm{real}}
SE(\hat\beta)
\end{lstlisting}
\resulthead
\[
GDP_{it}^{\mathrm{real}}
\qquad
SE(\hat\beta)
\]

\subsection{operatorname: named mathematical operators}

\sourcehead
\begin{lstlisting}
\operatorname{Cov}(X,Y)
\qquad
\operatorname{arg\,max}_{\theta} Q(\theta)
\end{lstlisting}
\resulthead
\[
\operatorname{Cov}(X,Y)
\qquad
\operatorname*{arg\,max}_{\theta} Q(\theta)
\]

For an operator you use repeatedly, define it once in the preamble:

\sourcehead
\begin{lstlisting}
\DeclareMathOperator{\plim}{plim}
\DeclareMathOperator*{\argmax}{arg\,max}
\end{lstlisting}

Then write \code{\textbackslash plim} or \code{\textbackslash argmax} throughout the paper.

\section{Math Alphabets: bold, blackboard, calligraphic, and Bold Greek}

These commands encode different mathematical roles:

\begin{center}
\begin{tabular}{lll}
\toprule
Source & Result & Typical use \\
\midrule
\code{\textbackslash mathbf\{x\}} & $\mathbf{x}$ & Latin vector/matrix \\
\code{\textbackslash bm\{\textbackslash beta\}} & $\bm{\beta}$ & bold Greek/vector \\
\code{\textbackslash mathbb\{R\}} & $\mathbb{R}$ & number sets \\
\code{\textbackslash mathcal\{F\}} & $\mathcal{F}$ & sets, sigma-fields, objective classes \\
\code{\textbackslash mathrm\{GDP\}} & $\mathrm{GDP}$ & upright abbreviation in math \\
\bottomrule
\end{tabular}
\end{center}

\sourcehead
\begin{lstlisting}
\mathbf{x}_i'\bm{\beta}
\qquad
\theta\in\mathbb{R}^K
\qquad
X\in\mathcal{F}
\end{lstlisting}
\resulthead
\[
\mathbf{x}_i'\bm{\beta}
\qquad
\theta\in\mathbb{R}^K
\qquad
X\in\mathcal{F}
\]

\warning{Do not use bold merely for emphasis inside equations. A notation style should be consistent: for example, bold lowercase for vectors and bold uppercase for matrices.}

\section{Delimiter Size: left/right and Manual Alternatives}

Ordinary parentheses do not grow automatically around tall material.

\sourcehead
\begin{lstlisting}
( \frac{A}{B} )
\qquad
\left( \frac{A}{B} \right)
\end{lstlisting}
\resulthead
\[
( \frac{A}{B} )
\qquad
\left( \frac{A}{B} \right)
\]

The same works with brackets, braces, absolute values, and norms:

\sourcehead
\begin{lstlisting}
\left[ \frac{A}{B} \right]
\left\{ \frac{A}{B} \right\}
\left| x_i \right|
\left\| \mathbf{x} \right\|
\end{lstlisting}

Sometimes automatic sizing is too large. Manual sizes give finer control:

\sourcehead
\begin{lstlisting}
\bigl( x \bigr)
\Bigl( x \Bigr)
\biggl( x \biggr)
\Biggl( x \Biggr)
\end{lstlisting}

\resulthead
\[
\bigl( x \bigr)
\qquad
\Bigl( x \Bigr)
\qquad
\biggl( x \biggr)
\qquad
\Biggl( x \Biggr)
\]

\subsection{One-sided delimiters with a dot}

\sourcehead
\begin{lstlisting}
\left. \frac{\partial f(x,\theta)}{\partial \theta}
\right|_{\theta=\theta_0}
\end{lstlisting}
\resulthead
\[
\left. \frac{\partial f(x,\theta)}{\partial \theta}
\right|_{\theta=\theta_0}
\]

The dot after \code{\textbackslash left} means ``no visible left delimiter.''

\section{Useful Annotation Commands: underbrace, overbrace, overset, and underset}

These commands can make a derivation easier to read when used sparingly.

\sourcehead
\begin{lstlisting}
\underbrace{\alpha + \beta x_i}_{\text{systematic component}}
+ \varepsilon_i
\end{lstlisting}
\resulthead
\[
\underbrace{\alpha + \beta x_i}_{\text{systematic component}}
+ \varepsilon_i
\]

\sourcehead
\begin{lstlisting}
\overbrace{x_1+\cdots+x_n}^{n\text{ observations}}
\end{lstlisting}
\resulthead
\[
\overbrace{x_1+\cdots+x_n}^{n\text{ observations}}
\]

Use \code{\textbackslash overset} and \code{\textbackslash underset} when you need a symbol directly above or below another relation:

\sourcehead
\begin{lstlisting}
\hat\theta_n \overset{p}{\longrightarrow} \theta_0
\qquad
x \underset{H_0}{\sim} F_0
\end{lstlisting}
\resulthead
\[
\hat\theta_n \overset{p}{\longrightarrow} \theta_0
\qquad
x \underset{H_0}{\sim} F_0
\]

\section{substack: Put Multiple Conditions Under a Sum, Max, or Limit}

A long subscript can be broken into multiple centered lines with \code{\textbackslash substack}.

\sourcehead
\begin{lstlisting}
\sum_{\substack{i=1,\ldots,n\\ t=1,\ldots,T}} y_{it}
\end{lstlisting}
\resulthead
\[
\sum_{\substack{i=1,\ldots,n\\ t=1,\ldots,T}} y_{it}
\]

Another common pattern is

\sourcehead
\begin{lstlisting}
\max_{\substack{c_1,c_2\\ c_1+c_2\le y}} u(c_1,c_2)
\end{lstlisting}
\resulthead
\[
\max_{\substack{c_1,c_2\\ c_1+c_2\le y}} u(c_1,c_2)
\]

\section{Dots: ldots, cdots, vdots, and ddots}

Different kinds of dots are conventional in different contexts.

\begin{center}
\begin{tabular}{lll}
\toprule
Command & Result & Typical use \\
\midrule
\code{\textbackslash ldots} & $\ldots$ & comma-separated lists: $1,\ldots,n$ \\
\code{\textbackslash cdots} & $\cdots$ & centered operations: $x_1+\cdots+x_n$ \\
\code{\textbackslash vdots} & $\vdots$ & vertical continuation in a matrix \\
\code{\textbackslash ddots} & $\ddots$ & diagonal continuation in a matrix \\
\bottomrule
\end{tabular}
\end{center}

\sourcehead
\begin{lstlisting}
1,\ldots,n
\qquad
x_1+\cdots+x_n
\end{lstlisting}
\resulthead
\[
1,\ldots,n
\qquad
x_1+\cdots+x_n
\]

\section{Relations and Definitions You Will See Often}

You do not need to memorize all mathematical symbols, but several recur in economics papers:

\begin{center}
\begin{tabular}{llll}
\toprule
Source & Result & Source & Result \\
\midrule
\code{\textbackslash approx} & $\approx$ & \code{\textbackslash sim} & $\sim$ \\
\code{\textbackslash propto} & $\propto$ & \code{\textbackslash equiv} & $\equiv$ \\
\code{\textbackslash neq} & $\neq$ & \code{\textbackslash leq} & $\leq$ \\
\code{\textbackslash geq} & $\geq$ & \code{\textbackslash in} & $\in$ \\
\code{\textbackslash subseteq} & $\subseteq$ & \code{\textbackslash perp} & $\perp$ \\
\code{\textbackslash coloneqq} & $\coloneqq$ & \code{\textbackslash implies} & $\implies$ \\
\bottomrule
\end{tabular}
\end{center}

\sourcehead
\begin{lstlisting}
\theta \in \Theta \subseteq \mathbb{R}^K,
\qquad
m(\theta) \coloneqq \E[g(W_i,\theta)].
\end{lstlisting}
\resulthead
\[
\theta \in \Theta \subseteq \mathbb{R}^K,
\qquad
m(\theta) \coloneqq \E[g(W_i,\theta)].
\]

\tip{\code{\textbackslash coloneqq} is especially useful when you mean ``is defined as,'' while \code{=} should usually mean equality. The command is supplied here by \code{mathtools}.}

\section{Reference Shortcuts Beyond ref and eqref}

You already know \code{\textbackslash ref} and \code{\textbackslash eqref}. A few related commands are worth recognizing.

\begin{center}
\begin{tabularx}{0.96\textwidth}{lX}
\toprule
Command & What it does \\
\midrule
\code{\textbackslash ref\{tab:main\}} & Prints only the number, such as 3. \\
\code{\textbackslash eqref\{eq:model\}} & Prints the equation number with parentheses, such as (2). \\
\code{\textbackslash pageref\{tab:main\}} & Prints the page on which the label appears. \\
\code{\textbackslash autoref\{fig:main\}} & With \code{hyperref}, prints the object type and number, such as ``Figure 2.'' \\
\code{\textbackslash nameref\{sec:data\}} & Prints the title of the labeled section, such as ``Data.'' \\
\bottomrule
\end{tabularx}
\end{center}

\sourcehead
\begin{lstlisting}
See Table~\ref{tab:main} on page~\pageref{tab:main}.
The construction is described in \nameref{sec:data}.
\end{lstlisting}

\tip{For economics papers, explicit prose such as \code{Table\textasciitilde\textbackslash ref\{...\}} is often clearer than relying on \code{\textbackslash autoref} everywhere. But it is useful to know that automatic object names exist.}

\section{Links, Email Addresses, URLs, and Literal Paths}

Because this handbook loads \code{hyperref}, several commands are available.

\subsection{href: custom clickable text}

\sourcehead
\begin{lstlisting}
\href{https://www.nber.org/}{NBER website}
\end{lstlisting}
\resulthead
\href{https://www.nber.org/}{NBER website}

\subsection{url: print a URL safely}

\sourcehead
\begin{lstlisting}
\url{https://www.example.com/data?id=10&year=2026}
\end{lstlisting}

\subsection{mailto: clickable email}

\sourcehead
\begin{lstlisting}
\href{mailto:name@university.edu}{name@university.edu}
\end{lstlisting}

\subsection{nolinkurl: print URL-like text without a link}

\sourcehead
\begin{lstlisting}
\nolinkurl{data/raw/county_panel_2026.csv}
\end{lstlisting}
\resulthead
\nolinkurl{data/raw/county_panel_2026.csv}

\tip{\code{\textbackslash url} and \code{\textbackslash nolinkurl} are safer than manually escaping every underscore in a long URL or path.}

\section{Font Size Commands and Local Scope}

LaTeX offers relative font-size declarations:

\begin{center}
\begin{tabular}{llllllll}
\toprule
\code{\textbackslash tiny} & \code{\textbackslash scriptsize} & \code{\textbackslash footnotesize} & \code{\textbackslash small} & \code{\textbackslash normalsize} & \code{\textbackslash large} & \code{\textbackslash Large} & \code{\textbackslash LARGE} \\
\bottomrule
\end{tabular}
\end{center}

Use braces to limit the effect:

\sourcehead
\begin{lstlisting}
Normal text.

{\small This paragraph is smaller, but only inside these braces.}

Normal text again.
\end{lstlisting}

A common table pattern is

\sourcehead
\begin{lstlisting}
\begin{table}[htbp]
\centering
\small
\begin{tabular}{lcc}
...
\end{tabular}
\end{table}
\end{lstlisting}

\warning{Do not shrink a table to unreadable size merely to force it onto one page. First shorten headings, remove unnecessary columns, split panels, or use landscape orientation.}

\section{Line and Page Control: Know the Difference}

\begin{center}
\begin{tabularx}{0.96\textwidth}{lX}
\toprule
Command & Meaning \\
\midrule
\code{\textbackslash\textbackslash} & End a line inside environments that expect line breaks; do not use it as a general paragraph command. \\
\code{\textbackslash newline} & Force a line break in ordinary text. Rarely needed in papers. \\
\code{\textbackslash linebreak} & Encourage/force a line break; optional strength \code{[0]}--\code{[4]}. \\
\code{\textbackslash newpage} & Start a new page immediately, but queued floats may remain pending. \\
\code{\textbackslash clearpage} & Print pending floats, then start a new page. \\
\code{\textbackslash pagebreak} & Encourage/force a page break; optional strength \code{[0]}--\code{[4]}. \\
\code{\textbackslash nopagebreak} & Discourage a page break at this point. \\
\code{\textbackslash FloatBarrier} & From \code{placeins}; stop earlier floats from moving beyond this point. \\
\bottomrule
\end{tabularx}
\end{center}

For an appendix that should begin after all main-text figures and tables have appeared, a common pattern is

\sourcehead
\begin{lstlisting}
\clearpage
\appendix
\section{Additional Results}
\end{lstlisting}

\section{Footnote Tricks: footnotemark and footnotetext}

Usually you should write

\sourcehead
\begin{lstlisting}
The data come from the administrative records.\footnote{Access requires approval.}
\end{lstlisting}

Sometimes the marker and the footnote text cannot be placed together, especially inside certain tables or captions. Then you may separate them:

\sourcehead
\begin{lstlisting}
Some text\footnotemark

...later...

\footnotetext{This is the corresponding footnote text.}
\end{lstlisting}

\warning{Footnotes inside tables and floats can be tricky. For economics regression tables, a table-note environment such as \code{threeparttable} is usually cleaner than ordinary footnotes.}

\section{Input, Include, and Includeonly}

For a real research paper, \code{\textbackslash input} is one of the most useful small commands.

\sourcehead
\begin{lstlisting}
\input{sections/introduction}
\input{sections/data}
\input{sections/results}
\input{tables/main_results}
\end{lstlisting}

LaTeX inserts each file exactly where the command appears. The inserted file usually does \emph{not} contain its own \code{\textbackslash documentclass} or \code{document} environment.

\code{\textbackslash include} is heavier: it starts the included material on a new page and creates separate auxiliary files. It is more common for books or dissertations than ordinary journal articles.

\sourcehead
\begin{lstlisting}
\include{chapter1}
\include{chapter2}
\end{lstlisting}

During a very large project, \code{\textbackslash includeonly} can compile only selected included files:

\sourcehead
\begin{lstlisting}
\includeonly{chapter2}
\end{lstlisting}

\tip{For an economics paper, default to \code{\textbackslash input}. It works naturally for sections, generated tables, appendix material, and reusable snippets.}

\section{newcommand and ensuremath: Make Your Own Shortcuts}

If a long term or notation appears repeatedly, define it once.

\sourcehead
\begin{lstlisting}
% in the preamble
\newcommand{\outcome}{labor-force participation}

% in the paper
Our main outcome is \outcome.
\end{lstlisting}

For a notation command that should work both inside and outside math mode, \code{\textbackslash ensuremath} is useful:

\sourcehead
\begin{lstlisting}
\newcommand{\ATE}{\ensuremath{\mathrm{ATE}}}

The parameter \ATE\ is reported in Table~\ref{tab:main}.
\end{lstlisting}

For commands that take arguments:

\sourcehead
\begin{lstlisting}
\newcommand{\Econd}[2]{\E\!\left[#1\mid #2\right]}

\Econd{Y_i}{X_i}
\end{lstlisting}
\resulthead
\[
\E\!\left[Y_i\mid X_i\right]
\]

\warning{Custom commands should improve consistency, not make your source unreadable. A collaborator should be able to infer what \code{\textbackslash ATE} means; an opaque command such as \code{\textbackslash zzzA} is not helpful.}

\section{A Few Useful Source-Code Controls}

\subsection{Percent: comment out the rest of a source line}

\sourcehead
\begin{lstlisting}
This prints. % This does not print.
\end{lstlisting}

A percent sign can also suppress an unwanted source-code newline after a command definition or line break. Beginners do not need this trick often, but you may see it in templates:

\sourcehead
\begin{lstlisting}
\newcommand{\myname}{%
  Ning Zhang%
}
\end{lstlisting}

\subsection{Temporarily hide a larger block}

A simple package-free method is

\sourcehead
\begin{lstlisting}
\iffalse
This entire block is ignored by LaTeX.
It may contain several paragraphs.
\fi
\end{lstlisting}

\warning{Do not leave large amounts of obsolete material hidden in a final project forever. Version control (for example, Git) is a better long-run archive than hundreds of commented-out lines.}

\section{Small Commands Worth Recognizing but Not Overusing}

These are useful when you encounter them in coauthors' files or journal templates.

\begin{center}
\begin{tabularx}{0.96\textwidth}{lX}
\toprule
Command & Purpose \\
\midrule
\code{\textbackslash raisebox\{1pt\}\{text\}} & Move a small object vertically. Mostly a fine-tuning tool. \\
\code{\textbackslash resizebox\{\textbackslash textwidth\}\{!\}\{...\}} & Scale an object to a width. Useful, but often abused on tables. \\
\code{\textbackslash scalebox\{0.9\}\{...\}} & Uniformly scale an object. \\
\code{\textbackslash rotatebox\{90\}\{...\}} & Rotate text or a small object. \\
\code{\textbackslash mbox\{...\}} & Put material in an unbreakable horizontal box. \\
\code{\textbackslash fbox\{...\}} & Draw a visible frame around material; useful for debugging or draft callouts. \\
\code{\textbackslash rule\{width\}\{height\}} & Draw a rule or create a precisely sized invisible/visible block. \\
\code{\textbackslash strut} & Insert an invisible standard-height object; sometimes helps equalize row heights. \\
\bottomrule
\end{tabularx}
\end{center}

\sourcehead
\begin{lstlisting}
\rotatebox{90}{Outcome}
\qquad
\fbox{Draft result}
\end{lstlisting}
\resulthead
\rotatebox{90}{Outcome}
\qquad
\fbox{Draft result}

\section{Fast Lookup: ``I Need to...''}

\begin{center}
\begin{tabularx}{0.98\textwidth}{Xl}
\toprule
I need to... & Try this first \\
\midrule
add a small math space & \code{\textbackslash,}, \code{\textbackslash:}, or \code{\textbackslash;} \\
add a larger math space & \code{\textbackslash quad} or \code{\textbackslash qquad} \\
tighten two math symbols slightly & \code{\textbackslash!} \\
push text to the right edge & \code{\textbackslash hfill} \\
add a precise horizontal gap & \code{\textbackslash hspace\{...\}} \\
add a little vertical separation & \code{\textbackslash smallskip}, \code{\textbackslash medskip}, \code{\textbackslash bigskip} \\
keep ``Table 3'' on one line & \code{Table\textasciitilde\textbackslash ref\{...\}} \\
remove paragraph indentation & \code{\textbackslash noindent} \\
center material in a figure/table & \code{\textbackslash centering} \\
align equations at the equals sign & put \code{\&} before \code{=} inside \code{align} \\
remove one equation number & \code{\textbackslash notag} \\
insert prose between aligned equations & \code{\textbackslash intertext\{...\}} \\
put ordinary words in math & \code{\textbackslash text\{...\}} \\
write an upright label in math & \code{\textbackslash mathrm\{...\}} \\
make Greek bold & \code{\textbackslash bm\{\textbackslash beta\}} \\
auto-size parentheses & \code{\textbackslash left(...\textbackslash right)} \\
label a component of an equation & \code{\textbackslash underbrace\{...\}\_\{...\}} \\
put two lines under a sum/max & \code{\textbackslash substack\{...\textbackslash\textbackslash...\}} \\
write a sequence of values & \code{1,\textbackslash ldots,n} \\
write a sum with centered dots & \code{x\_1+\textbackslash cdots+x\_n} \\
refer to a page & \code{\textbackslash pageref\{...\}} \\
print the name of a section & \code{\textbackslash nameref\{...\}} \\
insert a clickable link & \code{\textbackslash href\{URL\}\{text\}} \\
print a file path safely & \code{\textbackslash nolinkurl\{...\}} \\
start appendix after all floats & \code{\textbackslash clearpage} then \code{\textbackslash appendix} \\
insert another .tex file & \code{\textbackslash input\{...\}} \\
define a repeated abbreviation & \code{\textbackslash newcommand} \\
keep a command usable in/out of math & \code{\textbackslash ensuremath} \\
make a local font-size change & \code{\{\textbackslash small ...\}} \\
stop floats crossing a boundary & \code{\textbackslash FloatBarrier} \\
\bottomrule
\end{tabularx}
\end{center}

\tip{This section is deliberately repetitive. A beginner learns these commands by seeing them repeatedly in different contexts, not by memorizing a vocabulary list once.}



% ============================================================
\part{Gap Check: More Small Commands Beginners Commonly Miss}
% ============================================================

This part is a second-pass audit of the handbook. The commands here are not exotic tricks. They are small pieces of LaTeX that experienced users reach for repeatedly, especially when polishing economics papers, referee responses, appendices, and course notes.

\section{A Command Name Can Eat the Space After It}

A beginner often writes a command followed by ordinary text and is surprised that the space disappears.

\sourcehead
\begin{lstlisting}
\LaTeX is useful.
\LaTeX{} is useful.
\LaTeX\ is useful.
\end{lstlisting}

\resulthead
\LaTeX is useful.

\LaTeX{} is useful.

\LaTeX\ is useful.

\whyhead
A command made only of letters, such as \code{\textbackslash LaTeX}, consumes the following source-space while TeX decides where the command name ends. The empty braces in \code{\{\}} end the command without printing anything. The control-space \code{\textbackslash\ } explicitly inserts a normal space.

\tip{For commands that produce text, \code{\{\}} is usually the cleanest way to preserve the following space: \code{\textbackslash LaTeX\{\} paper}.}

\subsection{Abbreviations followed by a normal space}

LaTeX treats a period after a lowercase or uppercase sequence as potentially sentence-ending. When the period is part of an abbreviation, a control-space can be useful:

\sourcehead
\begin{lstlisting}
Prof.\ Smith
U.S.\ data
Fig.\ 2
\end{lstlisting}

\resulthead
Prof.\ Smith \qquad U.S.\ data \qquad Fig.\ 2

\section{More Spacing Tools: enspace, hspace*, and vfill}

The earlier spacing chapter covered the commands you will use most often. Three additional ones are worth recognizing.

\begin{center}
\begin{tabularx}{0.95\textwidth}{@{}lX@{}}
\toprule
Command & What it does \\
\midrule
\code{\textbackslash enspace} & Inserts a horizontal space roughly half an em wide. \\
\code{\textbackslash hspace*\{...\}} & Like \code{\textbackslash hspace}, but the space is preserved even at the beginning or end of a line. \\
\code{\textbackslash vfill} & Inserts flexible vertical space that expands to fill the available height. \\
\bottomrule
\end{tabularx}
\end{center}

\sourcehead
\begin{lstlisting}
$A\enspace B$

\hspace*{2em}Indented with explicit space.
\end{lstlisting}

\resulthead
$A\enspace B$

\hspace*{2em}Indented with explicit space.

A simple use of \code{\textbackslash vfill} is to push material toward opposite ends of a fixed-height box:

\resulthead
\begin{center}
\fbox{%
\begin{minipage}[t][1.2in][t]{0.62\textwidth}
Top line
\vfill
Bottom line
\end{minipage}}
\end{center}

\sourcehead
\begin{lstlisting}
\begin{minipage}[t][1.2in][t]{0.62\textwidth}
Top line
\vfill
Bottom line
\end{minipage}
\end{lstlisting}

\warning{Do not use manual spacing to repair the overall layout of a paper. If many paragraphs require custom \code{\textbackslash hspace} or \code{\textbackslash vspace}, the underlying structure probably needs fixing.}

\section{Lists: itemize, enumerate, and description}

Lists appear in theory assumptions, data descriptions, referee responses, research plans, and appendices.

\subsection{Bullets with itemize}

\sourcehead
\begin{lstlisting}
\begin{itemize}
  \item Administrative employment records
  \item County-level income data
  \item Firm characteristics
\end{itemize}
\end{lstlisting}

\resulthead
\begin{itemize}
  \item Administrative employment records
  \item County-level income data
  \item Firm characteristics
\end{itemize}

\subsection{Numbered lists with enumerate}

\sourcehead
\begin{lstlisting}
\begin{enumerate}
  \item Construct the analysis sample.
  \item Estimate the baseline specification.
  \item Report robustness checks.
\end{enumerate}
\end{lstlisting}

\resulthead
\begin{enumerate}
  \item Construct the analysis sample.
  \item Estimate the baseline specification.
  \item Report robustness checks.
\end{enumerate}

Because this handbook loads \code{enumitem}, you can change the label style easily:

\sourcehead
\begin{lstlisting}
\begin{enumerate}[label=(\roman*)]
  \item relevance
  \item exclusion
  \item monotonicity
\end{enumerate}
\end{lstlisting}

\resulthead
\begin{enumerate}[label=(\roman*)]
  \item relevance
  \item exclusion
  \item monotonicity
\end{enumerate}

\subsection{Term--definition lists with description}

\sourcehead
\begin{lstlisting}
\begin{description}
  \item[Outcome.] Annual earnings.
  \item[Treatment.] Program eligibility.
  \item[Sample.] Workers aged 25--54.
\end{description}
\end{lstlisting}

\resulthead
\begin{description}
  \item[Outcome.] Annual earnings.
  \item[Treatment.] Program eligibility.
  \item[Sample.] Workers aged 25--54.
\end{description}

\section{Fraction Size: frac, dfrac, tfrac, and displaystyle}

The ordinary command \code{\textbackslash frac} adapts to its context. Sometimes you want explicit control.

\begin{center}
\begin{tabularx}{0.95\textwidth}{@{}lX@{}}
\toprule
Command & Use \\
\midrule
\code{\textbackslash frac\{a\}\{b\}} & Normal context-sensitive fraction. \\
\code{\textbackslash dfrac\{a\}\{b\}} & Force display-style fraction. \\
\code{\textbackslash tfrac\{a\}\{b\}} & Force compact text-style fraction. \\
\code{\textbackslash displaystyle} & Force display-style math for the following expression. \\
\bottomrule
\end{tabularx}
\end{center}

\sourcehead
\begin{lstlisting}
Inline: $\frac{1}{n}\sum_i x_i$

Larger inline: $\dfrac{1}{n}\sum_i x_i$

Compact: $\tfrac{1}{2}\sigma^2$

$\displaystyle \sum_{i=1}^{n} x_i$
\end{lstlisting}

\resulthead
Inline: $\frac{1}{n}\sum_i x_i$

Larger inline: $\dfrac{1}{n}\sum_i x_i$

Compact: $\tfrac{1}{2}\sigma^2$

$\displaystyle \sum_{i=1}^{n} x_i$

\warning{Large display-style fractions can make a line of prose look uneven. Use \code{\textbackslash dfrac} only when the larger fraction genuinely improves readability.}

\section{Absolute Values, Norms, Angles, and Conditional Bars}

For mathematical delimiters, use semantic delimiter commands rather than typing keyboard characters blindly.

\sourcehead
\begin{lstlisting}
\lvert x \rvert
\lVert \beta \rVert
\langle x,y \rangle
\left\lvert \frac{x}{y} \right\rvert
\left\{x \middle| x>0\right\}
\end{lstlisting}

\resulthead
\[
\lvert x \rvert
\qquad
\lVert \beta \rVert
\qquad
\langle x,y \rangle
\qquad
\left\lvert \frac{x}{y} \right\rvert
\qquad
\left\{x \middle| x>0\right\}.
\]

\whyhead
\begin{itemize}
  \item \code{\textbackslash lvert...\textbackslash rvert}: absolute values.
  \item \code{\textbackslash lVert...\textbackslash rVert}: norms.
  \item \code{\textbackslash langle...\textbackslash rangle}: angle brackets or inner products.
  \item \code{\textbackslash middle|}: a conditional bar whose height matches surrounding \code{\textbackslash left} and \code{\textbackslash right} delimiters.
\end{itemize}

\section{Binomial Coefficients, Boxes, and Wider Accents}

A few compact commands are useful in theory notes, proofs, and appendices.

\sourcehead
\begin{lstlisting}
\binom{n}{k}
\boxed{\beta>0}
\widehat{\beta_1+\beta_2}
\widetilde{Y}_{it}
\overline{X}_g
\vec{x}
\dot{x}_t
\ddot{x}_t
\end{lstlisting}

\resulthead
\[
\binom{n}{k}
\qquad
\boxed{\beta>0}
\qquad
\widehat{\beta_1+\beta_2}
\qquad
\widetilde{Y}_{it}
\qquad
\overline{X}_g
\qquad
\vec{x}
\qquad
\dot{x}_t
\qquad
\ddot{x}_t.
\]

\tip{\code{\textbackslash widehat} and \code{\textbackslash widetilde} are preferable to \code{\textbackslash hat} and \code{\textbackslash tilde} when the accent must span several symbols.}

\section{Two More Equation Environments: gather and multline}

You already learned \code{equation}, \code{align}, \code{split}, and \code{aligned}. Two other \code{amsmath} environments solve different layout problems.

\subsection{gather: centered equations with no alignment point}

Use \code{gather} when several displayed equations belong together but do not need their equals signs aligned.

\sourcehead
\begin{lstlisting}
\begin{gather*}
E[u_i]=0,\\
\operatorname{Var}(u_i)=\sigma^2,\\
E[u_i x_i]=0.
\end{gather*}
\end{lstlisting}

\resulthead
\begin{gather*}
E[u_i]=0,\\
\operatorname{Var}(u_i)=\sigma^2,\\
E[u_i x_i]=0.
\end{gather*}

\subsection{multline: one long equation over several lines}

Use \code{multline} for one long equation that does not have a natural alignment point.

\sourcehead
\begin{lstlisting}
\begin{multline*}
Q(\theta)
= \sum_{i=1}^{n} q_i(\theta)
  + \lambda_1 P_1(\theta) \\
  + \lambda_2 P_2(\theta)
  + \lambda_3 P_3(\theta).
\end{multline*}
\end{lstlisting}

\resulthead
\begin{multline*}
Q(\theta)
= \sum_{i=1}^{n} q_i(\theta)
  + \lambda_1 P_1(\theta) \\
  + \lambda_2 P_2(\theta)
  + \lambda_3 P_3(\theta).
\end{multline*}

\subsection{tag: give an equation a custom visible tag}

\sourcehead
\begin{lstlisting}
\begin{equation*}
Y_i = \alpha + \beta X_i + u_i.
\tag{A1}
\end{equation*}
\end{lstlisting}

\resulthead
\begin{equation*}
Y_i = \alpha + \beta X_i + u_i.
\tag{A1}
\end{equation*}

\warning{Do not manually tag ordinary equations just to obtain the number you want. Use automatic numbering and \code{\textbackslash label}/\code{\textbackslash ref}. Manual tags are mainly useful when reproducing an external numbering convention.}

\section{Allowing a Break in a Long Mathematical Expression}

A long inline expression may create an overfull line because TeX has too few legal break points. \code{\textbackslash allowbreak} adds a possible break without forcing one.

\sourcehead
\begin{lstlisting}
$X_{it}=(age_{it},\allowbreak education_{it},
\allowbreak experience_{it},\allowbreak income_{it})$
\end{lstlisting}

\whyhead
If the line fits, nothing visible happens. If the line is too long, LaTeX may break at an \code{\textbackslash allowbreak}. This is useful for long lists of variables or parameters.

\warning{First ask whether the expression should be displayed on its own line. Do not sprinkle \code{\textbackslash allowbreak} everywhere to hide poor equation design.}

\section{textwidth, linewidth, and columnwidth: Which Width Should You Use?}

These length commands look similar but refer to different available widths.

\begin{center}
\begin{tabularx}{0.96\textwidth}{@{}lX@{}}
\toprule
Length & Meaning \\
\midrule
\code{\textbackslash textwidth} & Width of the main text block on the page. \\
\code{\textbackslash linewidth} & Width currently available in the local environment. Inside a \code{minipage} or \code{subfigure}, this is usually the safest choice. \\
\code{\textbackslash columnwidth} & Width of one column. In a one-column article it is usually the same as \code{\textbackslash textwidth}; in a two-column document it is smaller. \\
\bottomrule
\end{tabularx}
\end{center}

A common figure pattern is:

\sourcehead
\begin{lstlisting}
\begin{figure}
  \centering
  \includegraphics[width=0.8\textwidth]{main.pdf}
\end{figure}

\begin{subfigure}{0.48\textwidth}
  \includegraphics[width=\linewidth]{panel_a.pdf}
\end{subfigure}
\end{lstlisting}

\tip{Inside a \code{subfigure}, \code{minipage}, table cell, or other nested box, \code{\textbackslash linewidth} often means what you actually want.}

\section{graphicspath: Stop Repeating the Figures Folder}

If nearly all graphics live in the same folder, tell \code{graphicx} where to look.

\sourcehead
\begin{lstlisting}
\graphicspath{{figures/}{figures/appendix/}}

% Later:
\includegraphics[width=.8\textwidth]{main_result.pdf}
\end{lstlisting}

\whyhead
LaTeX first looks in the current directory, then in the listed graphics directories. The double braces are intentional: each path is wrapped in its own pair of braces.

\tip{Many research projects still prefer explicit paths such as \code{figures/main\_result.pdf} because they make the source easier to audit. Either convention is fine if the project is consistent.}

\section{Four Table Controls That Save a Lot of Frustration}

\subsection{tabcolsep: horizontal padding between columns}

\sourcehead
\begin{lstlisting}
{\setlength{\tabcolsep}{4pt}
\begin{tabular}{lcc}
...
\end{tabular}}
\end{lstlisting}

The braces keep the change local to that table.

\subsection{arraystretch: vertical row spacing}

\sourcehead
\begin{lstlisting}
{\renewcommand{\arraystretch}{1.15}
\begin{tabular}{lcc}
...
\end{tabular}}
\end{lstlisting}

A value larger than 1 makes rows taller.

\subsection{p\{width\}: a column that wraps long text}

\sourcehead
\begin{lstlisting}
\begin{tabular}{p{0.25\textwidth}p{0.65\textwidth}}
Variable & Definition \\
...
\end{tabular}
\end{lstlisting}

Unlike \code{l}, \code{c}, and \code{r}, a \code{p\{...\}} column wraps paragraphs automatically.

\subsection{@\{\}: remove or customize edge padding}

\sourcehead
\begin{lstlisting}
\begin{tabular}{@{}lcc@{}}
...
\end{tabular}
\end{lstlisting}

The empty \code{@\{\}} at both ends removes the normal outer column padding. This is one reason professional tables often align more cleanly with surrounding text.

\subsection{shortstack: a quick two-line column heading}

\sourcehead
\begin{lstlisting}
\begin{tabular}{lcc}
\toprule
 & \shortstack{Annual\\earnings}
 & \shortstack{Employment\\rate} \\
\midrule
Treatment & 0.12 & 0.03 \\
\bottomrule
\end{tabular}
\end{lstlisting}

\resulthead
\begin{center}
\begin{tabular}{lcc}
\toprule
 & \shortstack{Annual\\earnings}
 & \shortstack{Employment\\rate} \\
\midrule
Treatment & 0.12 & 0.03 \\
\bottomrule
\end{tabular}
\end{center}

\warning{If headings are becoming very complicated, rethink the table structure before stacking three or four lines into every cell.}

\section{minipage and parbox: Small Independent Text Blocks}

Both commands create a box with its own line width.

\subsection{minipage: best for multi-line blocks}

\sourcehead
\begin{lstlisting}
\begin{minipage}[t]{0.46\textwidth}
\textbf{Sample A}

Workers aged 25--40.
\end{minipage}
\hfill
\begin{minipage}[t]{0.46\textwidth}
\textbf{Sample B}

Workers aged 41--54.
\end{minipage}
\end{lstlisting}

\resulthead
\begin{minipage}[t]{0.46\textwidth}
\textbf{Sample A}

Workers aged 25--40.
\end{minipage}
\hfill
\begin{minipage}[t]{0.46\textwidth}
\textbf{Sample B}

Workers aged 41--54.
\end{minipage}

\subsection{parbox: convenient for a smaller inline text box}

\sourcehead
\begin{lstlisting}
\parbox{0.65\textwidth}{
This note wraps within a fixed-width box.
}
\end{lstlisting}

\resulthead
\parbox{0.65\textwidth}{
This note wraps within a fixed-width box.
}

\section{More natbib Citation Commands Worth Knowing}

The first citation chapter covered \code{\textbackslash citet} and \code{\textbackslash citep}. These additional commands are useful when you need the pieces of a citation separately.

\begin{center}
\begin{tabularx}{0.96\textwidth}{@{}lXX@{}}
\toprule
Command & Meaning & Typical output \\
\midrule
\code{\textbackslash citeauthor\{key\}} & Author name only & Author \\
\code{\textbackslash citeyear\{key\}} & Year only & 2020 \\
\code{\textbackslash citealp\{key\}} & Author-year without outer parentheses & Author, 2020 \\
\code{\textbackslash citealt\{key\}} & Textual-style citation without parentheses around year & Author 2020 \\
\code{\textbackslash nocite\{key\}} & Put an item in the bibliography without citing it in the text & no visible text here \\
\code{\textbackslash nocite\{*\}} & Include every entry from the bibliography database & no visible text here \\
\bottomrule
\end{tabularx}
\end{center}

\sourcehead
\begin{lstlisting}
\citeauthor{acemoglu2024}
\citeyear{acemoglu2024}
\citealp{acemoglu2024}
\nocite{workingpaper2025}
\end{lstlisting}

\warning{Do not manually type author names and years merely to imitate a citation. If the statement is a citation, let the bibliography system generate it.}

\section{Page Numbering and Page Style Commands}

Journal templates usually control these settings for you, but they are useful for working papers, appendices, referee responses, and course documents.

\begin{center}
\begin{tabularx}{0.96\textwidth}{@{}lX@{}}
\toprule
Command & Purpose \\
\midrule
\code{\textbackslash pagestyle\{plain\}} & Use a simple page style, usually page numbers. \\
\code{\textbackslash pagestyle\{empty\}} & Suppress headers and footers on all following pages. \\
\code{\textbackslash thispagestyle\{empty\}} & Suppress header/footer on this page only. \\
\code{\textbackslash pagenumbering\{roman\}} & Use i, ii, iii, ... and reset the page counter. \\
\code{\textbackslash pagenumbering\{arabic\}} & Use 1, 2, 3, ... and reset the page counter. \\
\code{\textbackslash setcounter\{page\}\{1\}} & Set the current page counter manually. \\
\bottomrule
\end{tabularx}
\end{center}

\sourcehead
\begin{lstlisting}
\thispagestyle{empty}

\pagenumbering{roman}
% front matter

\clearpage
\pagenumbering{arabic}
% main paper
\end{lstlisting}

\warning{Do not override page styles in a journal or conference template unless the instructions explicitly allow it.}

\section{Unnumbered Sections That Still Appear in the Contents}

The starred form removes automatic numbering:

\sourcehead
\begin{lstlisting}
\section*{Acknowledgments}
\end{lstlisting}

If you also want the unnumbered section in a table of contents:

\sourcehead
\begin{lstlisting}
\section*{Acknowledgments}
\addcontentsline{toc}{section}{Acknowledgments}
\end{lstlisting}

\whyhead
The three arguments of \code{\textbackslash addcontentsline} mean: write to the table of contents (\code{toc}), use section-level formatting (\code{section}), and print the text \code{Acknowledgments}.

\tip{Most journal articles do not need a table of contents. This command is more relevant for dissertations, long reports, notes, and handbooks like this one.}

\section{Line Breaking and Hyphenation: Small Rescue Tools}

\subsection{newline versus a new paragraph}

\sourcehead
\begin{lstlisting}
First line.\newline
Second line.

New paragraph after a blank source line.
\end{lstlisting}

\resulthead
First line.\newline
Second line.

New paragraph after a blank source line.

\whyhead
\code{\textbackslash newline} starts a new line but does not start a new paragraph. In ordinary prose, a blank line is usually the correct way to begin a new paragraph.

\subsection{A discretionary hyphen with \textbackslash-}

If LaTeX refuses to break a rare or technical word well, you can mark an allowed hyphenation point:

\sourcehead
\begin{lstlisting}
hetero\-geneity
counter\-factual
\end{lstlisting}

The hyphen appears only if LaTeX actually breaks the word at that point.

\subsection{Declare hyphenation for a whole document}

\sourcehead
\begin{lstlisting}
\hyphenation{hetero-geneity counter-factual}
\end{lstlisting}

Put \code{\textbackslash hyphenation} in the preamble.

\subsection{emergencystretch: a last-resort paragraph safety valve}

\sourcehead
\begin{lstlisting}
\setlength{\emergencystretch}{2em}
\end{lstlisting}

This gives TeX a little extra flexibility when it is struggling to avoid overfull lines.

\warning{Do not use \code{\textbackslash emergencystretch} as a substitute for fixing a genuinely unbreakable URL, a too-wide equation, or a bad table. It is a global typography adjustment.}

\section{One More Useful Distinction: center vs. centering}

Both can center material, but they behave differently.

\sourcehead
\begin{lstlisting}
\begin{center}
Centered paragraph.
\end{center}

\begin{figure}
  \centering
  \includegraphics{figure.pdf}
\end{figure}
\end{lstlisting}

\whyhead
The \code{center} environment adds vertical space before and after itself. The declaration \code{\textbackslash centering} changes alignment without adding that extra environment spacing. Inside \code{figure} and \code{table}, \code{\textbackslash centering} is normally preferred.

\section{A Second Fast Lookup: Commands Added After the Gap Check}

\begin{center}
\begin{tabularx}{0.98\textwidth}{@{}Xl@{}}
\toprule
I need to... & Try this \\
\midrule
keep a space after a text-producing command & \code{\{\}} or \code{\textbackslash\ } \\
insert a medium fixed horizontal space & \code{\textbackslash enspace} \\
preserve horizontal space at a line edge & \code{\textbackslash hspace*} \\
push material vertically apart & \code{\textbackslash vfill} \\
make bullets & \code{itemize} \\
make numbered conditions & \code{enumerate} \\
make term--definition entries & \code{description} \\
force a large fraction & \code{\textbackslash dfrac} \\
force a compact fraction & \code{\textbackslash tfrac} \\
write an absolute value & \code{\textbackslash lvert...\textbackslash rvert} \\
write a norm & \code{\textbackslash lVert...\textbackslash rVert} \\
write an inner product & \code{\textbackslash langle...\textbackslash rangle} \\
write a binomial coefficient & \code{\textbackslash binom} \\
span an accent over several symbols & \code{\textbackslash widehat}, \code{\textbackslash widetilde} \\
center several unrelated equations & \code{gather} \\
break one long displayed equation & \code{multline} \\
offer a legal break inside long inline math & \code{\textbackslash allowbreak} \\
size a graphic to its local container & \code{\textbackslash linewidth} \\
declare default figure folders & \code{\textbackslash graphicspath} \\
tighten table columns & \code{\textbackslash tabcolsep} \\
increase table row height & \code{\textbackslash arraystretch} \\
wrap long text in a table column & \code{p\{...\}} \\
make a two-line table header quickly & \code{\textbackslash shortstack} \\
put two text blocks side by side & \code{minipage} \\
print only a citation author & \code{\textbackslash citeauthor} \\
print only a citation year & \code{\textbackslash citeyear} \\
suppress the page style on one page & \code{\textbackslash thispagestyle\{empty\}} \\
add an unnumbered section to the TOC & \code{\textbackslash addcontentsline} \\
allow a manual word-break point & \code{\textbackslash-} \\
\bottomrule
\end{tabularx}
\end{center}

\tip{The goal is not to memorize this table. It gives you search vocabulary. When you know the name of the tool, you can quickly look up the exact syntax later.}



% ============================================================
\part{Showing Code and Referencing Things in Your Paper}
% ============================================================

\section{Three Different Jobs: Inline Code, Code Blocks, and References}

These three ideas look similar to a beginner, but they do different jobs:

\begin{center}
\begin{tabularx}{0.94\textwidth}{@{}p{0.22\textwidth}X@{}}
\toprule
\textbf{What you want} & \textbf{Typical LaTeX tool} \\
\midrule
Show a command literally inside a sentence & \code{\textbackslash verb}, \code{\textbackslash lstinline}, or \code{\textbackslash texttt} \\
Show several lines of source code & \code{verbatim} or \code{lstlisting} \\
Refer to a figure, table, equation, or section & \code{\textbackslash label} + \code{\textbackslash ref}, \code{\textbackslash autoref}, or \code{\textbackslash cref} \\
\bottomrule
\end{tabularx}
\end{center}

\tip{When you want the PDF to literally show a command such as \texttt{\textbackslash qquad}, you are \emph{typesetting code}. When you write \code{Figure\string~\textbackslash ref\{fig:main\}}, you are creating a \emph{cross-reference}.}

\section{Inline Code: Showing a LaTeX Command Literally}

Suppose you want your PDF to say: ``Use \verb|\qquad| for a large math-space.'' If you simply type \code{\textbackslash qquad}, LaTeX tries to execute the command. To show the command itself, use a verbatim-style command.

\subsection{The simplest tool: verb}

\sourcehead
\begin{lstlisting}
Use \verb|\qquad| for a large horizontal space.
Use \verb|\textwidth| for the page text width.
\end{lstlisting}

\resulthead
Use \verb|\qquad| for a large horizontal space.
Use \verb|\textwidth| for the page text width.

\whyhead
The character immediately after \code{\textbackslash verb} is a delimiter. Here the delimiter is \code{|}. LaTeX prints everything between the two delimiters literally.

You can choose another delimiter when the text itself contains \code{|}:

\begin{lstlisting}
\verb+\qquad+
\verb!data_file.dta!
\end{lstlisting}

\warning{Do not write braces after \code{\textbackslash verb}. It is \code{\textbackslash verb|text|}, not \code{\textbackslash verb\{text\}}. Also avoid \code{\textbackslash verb} inside section titles and captions.}

\subsection{lstinline: inline code that matches your listings style}

Because this handbook already uses the \code{listings} package, you can also write:

\sourcehead
\begin{lstlisting}
The command \lstinline|\qquad| inserts a large space.
The file \lstinline|analysis.R| contains the main estimation code.
\end{lstlisting}

\resulthead
The command \lstinline|\qquad| inserts a large space.
The file \lstinline|analysis.R| contains the main estimation code.

\tip{For a paper about economics, you usually do not need inline code often. It is most useful in notes, teaching material, appendices, replication documentation, and a LaTeX handbook like this one.}

\subsection{texttt: typewriter font, but not truly verbatim}

\sourcehead
\begin{lstlisting}
\texttt{analysis.R}
\texttt{data\_clean.dta}
\end{lstlisting}

\resulthead
\texttt{analysis.R}\qquad \texttt{data\_clean.dta}

\code{\textbackslash texttt} changes the font, but LaTeX still interprets reserved characters. That is why the underscore must be escaped in \code{data\_clean.dta}. For arbitrary source code, \code{\textbackslash verb} or \code{\textbackslash lstinline} is easier.

\section{Code Blocks: verbatim vs. lstlisting}

\subsection{verbatim: the basic built-in code block}

The \code{verbatim} environment prints characters exactly as typed. Reserved characters such as \code{\%}, \code{\_}, \code{\&}, and \code{\textbackslash} do not need escaping inside it.

\sourcehead
\begin{lstlisting}
\begin{verbatim}
reg y x1 x2, robust
save results_main.dta, replace
\end{verbatim}
\end{lstlisting}

\resulthead
\begin{verbatim}
reg y x1 x2, robust
save results_main.dta, replace
\end{verbatim}

\whyhead
\code{verbatim} is simple and dependable. It is good when you only need to show literal text and do not care about captions, labels, frames, syntax highlighting, or line numbers.

\subsection{lstlisting: better for polished code blocks}

The \code{listings} package gives you much more control.

\sourcehead
\begin{verbatim}
\begin{lstlisting}[language=R]
df <- read.csv("data/analysis.csv")
model <- lm(y ~ x1 + x2, data = df)
summary(model)
\end{lstlisting}
\end{verbatim}

\resulthead
\begin{lstlisting}[language=R]
df <- read.csv("data/analysis.csv")
model <- lm(y ~ x1 + x2, data = df)
summary(model)
\end{lstlisting}

\tip{For Stata code, you can leave \code{language=} unspecified unless you define a custom Stata language. The code block will still print correctly.}

\section{Useful listings Options}

A code block can have a caption, line numbers, a label, and a language.

\sourcehead
\begin{verbatim}
\begin{lstlisting}[
  language=R,
  numbers=left,
  caption={Minimal data-cleaning example},
  label={lst:cleaning}
]
df <- subset(df, age >= 18)
df$log_income <- log(df$income)
\end{lstlisting}
\end{verbatim}

\resulthead
\begin{lstlisting}[
  language=R,
  numbers=left,
  caption={Minimal data-cleaning example},
  label={lst:cleaning-demo}
]
df <- subset(df, age >= 18)
df$log_income <- log(df$income)
\end{lstlisting}

Now you can refer to the listing just like a figure or table:

\sourcehead
\begin{lstlisting}
Listing~\ref{lst:cleaning-demo} shows the cleaning step.
\end{lstlisting}

\resulthead
Listing~\ref{lst:cleaning-demo} shows the cleaning step.

\section{Include Code from an External File}

In a research project, it is often better not to copy code into the paper manually. The \code{listings} package can read a source file directly.

\sourcehead
\begin{lstlisting}
\lstinputlisting[language=R]{code/main_analysis.R}
\end{lstlisting}

You can also include only part of a file:

\begin{lstlisting}
\lstinputlisting[
  language=Python,
  firstline=20,
  lastline=45
]{code/figure1.py}
\end{lstlisting}

\tip{This is useful for appendices, replication notes, teaching notes, or presentations. Main economics journal articles usually do not print large blocks of analysis code in the paper itself.}

\section{The Basic Figure Reference Workflow}

A figure has four jobs: create a float, insert the image, provide a caption, and attach a label.

\sourcehead
\begin{lstlisting}
\begin{figure}[htbp]
  \centering
  \includegraphics[width=0.75\linewidth]{figures/main_result.pdf}
  \caption{Main result}
  \label{fig:main-result}
\end{figure}
\end{lstlisting}

The most important ordering rule is:

\begin{center}
\textbf{caption first, label second}
\end{center}

Then refer to it in your prose:

\sourcehead
\begin{lstlisting}
Figure~\ref{fig:main-result} reports the main result.
\end{lstlisting}

If the figure is numbered 2, the PDF automatically says:

\resulthead
\textbf{Figure 2 reports the main result.}

\warning{Do not manually type ``Figure 2'' throughout your source. If the figure order changes, the number becomes wrong. Use \code{\textbackslash ref}.}

\section{ref, eqref, autoref, cref, and Cref}

There are several levels of convenience.

\begin{center}
\begin{tabularx}{0.96\textwidth}{@{}p{0.19\textwidth}p{0.27\textwidth}X@{}}
\toprule
\textbf{Command} & \textbf{Typical source} & \textbf{What it does} \\
\midrule
\code{\textbackslash ref} & \code{Figure\string~\textbackslash ref\{fig:main\}} & Prints only the number; you type ``Figure'' yourself. \\
\code{\textbackslash eqref} & \code{\textbackslash eqref\{eq:main\}} & Prints the equation number with parentheses. \\
\code{\textbackslash autoref} & \code{\textbackslash autoref\{fig:main\}} & \code{hyperref} automatically adds the object name. \\
\code{\textbackslash cref} & \code{\textbackslash cref\{fig:main\}} & \code{cleveref} automatically chooses ``figure'', ``table'', ``section'', etc. \\
\code{\textbackslash Cref} & \code{\textbackslash Cref\{fig:main\}} & Same as \code{\textbackslash cref}, but capitalized for the start of a sentence. \\
\bottomrule
\end{tabularx}
\end{center}

For example, suppose a figure has \code{\textbackslash label\{fig:demo-ref\}} and is numbered 3.

\sourcehead
\begin{lstlisting}
Figure~\ref{fig:demo-ref}
\autoref{fig:demo-ref}
\cref{fig:demo-ref}
\Cref{fig:demo-ref}
\end{lstlisting}

Conceptually, these produce:

\begin{quote}
Figure 3\\
Figure 3\\
figure 3\\
Figure 3
\end{quote}

\tip{For beginners, \code{Figure\string~\textbackslash ref\{...\}} is completely fine. Once your paper has many references, \code{cleveref} can reduce typing and prevent mistakes such as writing ``Table'' before a figure reference.}

\section{Referencing Several Objects at Once with cleveref}

One advantage of \code{cleveref} is that it can format several references together.

\sourcehead
\begin{lstlisting}
\cref{fig:main,fig:robustness}
\cref{tab:summary,tab:main,tab:robust}
\end{lstlisting}

Depending on the numbering, LaTeX can automatically produce text such as:

\begin{quote}
figures 2 and 3\\
tables 1--3
\end{quote}

You do not have to manually decide whether to write ``Figure'' or ``Figures.''

\section{Panel References in Multi-Panel Figures}

With the \code{subcaption} package, each panel can have its own label.

\sourcehead
\begin{lstlisting}
\begin{figure}[htbp]
  \centering
  \begin{subfigure}{0.47\textwidth}
    \includegraphics[width=\linewidth]{figures/panel_a.pdf}
    \caption{Full sample}
    \label{fig:panel-a-ref}
  \end{subfigure}
  \hfill
  \begin{subfigure}{0.47\textwidth}
    \includegraphics[width=\linewidth]{figures/panel_b.pdf}
    \caption{Restricted sample}
    \label{fig:panel-b-ref}
  \end{subfigure}
  \caption{Main results by sample}
  \label{fig:two-panel-ref}
\end{figure}
\end{lstlisting}

Then you can write:

\begin{lstlisting}
Figure~\ref{fig:two-panel-ref}\subref{fig:panel-a-ref}
Figure~\ref{fig:two-panel-ref}\subref{fig:panel-b-ref}
\end{lstlisting}

which gives forms such as \textbf{Figure 4(a)} and \textbf{Figure 4(b)}.

\section{The Same Reference System Works for Tables, Equations, and Sections}

\sourcehead
\begin{lstlisting}
\section{Data}
\label{sec:data-ref}

\begin{equation}
  y_i = \alpha + \beta x_i + \varepsilon_i.
  \label{eq:demo-ref}
\end{equation}

\begin{table}[htbp]
  \centering
  \caption{Summary statistics}
  \label{tab:summary-ref}
  ...
\end{table}
\end{lstlisting}

Then write:

\begin{lstlisting}
Section~\ref{sec:data-ref}
Equation~\eqref{eq:demo-ref}
Table~\ref{tab:summary-ref}
\end{lstlisting}

The pattern is always the same:

\begin{center}
\boxed{\texttt{create object} \;\longrightarrow\; \texttt{label it} \;\longrightarrow\; \texttt{reference the label}}
\end{center}

\section{A Label-Naming Habit That Scales to a Dissertation}

Use a short prefix that tells you what kind of object the label belongs to:

\begin{center}
\begin{tabular}{@{}ll@{}}
\toprule
\textbf{Prefix} & \textbf{Example} \\
\midrule
Section & \code{sec:data} \\
Equation & \code{eq:baseline} \\
Figure & \code{fig:event-study} \\
Table & \code{tab:main-results} \\
Appendix & \code{app:proofs} \\
Proposition & \code{prop:existence} \\
Assumption & \code{assump:sampling} \\
Code listing & \code{lst:cleaning} \\
\bottomrule
\end{tabular}
\end{center}

This becomes especially valuable when your editor offers autocomplete after you type \code{\textbackslash ref\{}.

\section{Common Reference Mistakes}

\begin{center}
\begin{tabularx}{0.96\textwidth}{@{}p{0.31\textwidth}X@{}}
\toprule
\textbf{Problem} & \textbf{What to check} \\
\midrule
Reference prints \textbf{??} & Compile again; then check that the label exists and is spelled identically. \\
Figure reference gives the wrong number & Put \code{\textbackslash label} after \code{\textbackslash caption}. \\
Two objects have the same reference number unexpectedly & Check for duplicated label names. Every label should be unique. \\
You typed Figure 3 manually and it became wrong & Replace hard-coded numbers with \code{\textbackslash ref}. \\
\code{\textbackslash ref} gives only ``3'', not ``Figure 3'' & That is normal. Type \code{Figure\string~\textbackslash ref\{...\}} or use \code{\textbackslash autoref}/\code{\textbackslash cref}. \\
Citation or cross-reference is stale & Usually two clean recompilations are enough. \\
\bottomrule
\end{tabularx}
\end{center}

\section{Fast Lookup: Code and References}

\begin{center}
\begin{tabularx}{0.96\textwidth}{@{}p{0.46\textwidth}X@{}}
\toprule
\textbf{I need to...} & \textbf{Use} \\
\midrule
Show \texttt{\textbackslash qquad} literally in a sentence & \code{\textbackslash verb|\textbackslash qquad|} or \code{\textbackslash lstinline|\textbackslash qquad|} \\
Show a filename in typewriter font & \code{\textbackslash texttt\{analysis.R\}} \\
Show several lines exactly as typed & \code{verbatim} \\
Show a polished code block & \code{lstlisting} \\
Number and reference a code block & \code{caption=} + \code{label=} in \code{lstlisting} \\
Read code directly from a file & \code{\textbackslash lstinputlisting\{...\}} \\
Insert an image & \code{\textbackslash includegraphics} \\
Give a figure a reference name & \code{\textbackslash label\{fig:...\}} after \code{\textbackslash caption} \\
Refer to the figure number & \code{Figure\string~\textbackslash ref\{fig:...\}} \\
Automatically print the object type & \code{\textbackslash autoref} or \code{\textbackslash cref} \\
Refer to a displayed equation & \code{\textbackslash eqref\{eq:...\}} \\
Refer to a figure panel & \code{\textbackslash subref\{fig:panel-a\}} \\
Reference several objects together & \code{\textbackslash cref\{fig:a,fig:b\}} \\
\bottomrule
\end{tabularx}
\end{center}


% ============================================================
\part{Tables: From First Table to Paper-Ready Table}
% ============================================================

\section{The Difference Between table and tabular}

This distinction confuses many beginners.

\begin{itemize}
  \item \code{tabular} creates the rows and columns.
  \item \code{table} is a floating container that gives you a caption, label, and flexible placement.
\end{itemize}

\section{Your First Table}

\sourcehead
\begin{lstlisting}
\begin{tabular}{lcc}
Variable & Mean & SD \\
Income   & 50.2 & 12.4 \\
Age      & 39.7 & 11.3 \\
\end{tabular}
\end{lstlisting}

\resulthead
\begin{center}
\begin{tabular}{lcc}
Variable & Mean & SD \\
Income   & 50.2 & 12.4 \\
Age      & 39.7 & 11.3 \\
\end{tabular}
\end{center}

The column declaration \code{\{lcc\}} means:

\begin{itemize}
  \item first column left aligned (\code{l});
  \item second centered (\code{c});
  \item third centered (\code{c}).
\end{itemize}

Every \code{\&} moves to the next column. Every \code{\textbackslash\textbackslash} ends the row.

\section{A Paper-Ready Table with booktabs}

\sourcehead
\begin{lstlisting}
\begin{table}[htbp]
\centering
\caption{Summary Statistics}
\label{tab:summary}
\begin{tabular}{lccc}
\toprule
Variable & Mean & SD & N \\
\midrule
Outcome  & 4.82 & 1.31 & 1,000 \\
Control  & 2.14 & 0.87 & 1,000 \\
\bottomrule
\end{tabular}
\end{table}
\end{lstlisting}

\resulthead
\begin{table}[H]
\centering
\caption{Summary Statistics}
\label{tab:summary-demo}
\begin{tabular}{lccc}
\toprule
Variable & Mean & SD & N \\
\midrule
Outcome  & 4.82 & 1.31 & 1,000 \\
Control  & 2.14 & 0.87 & 1,000 \\
\bottomrule
\end{tabular}
\end{table}

Use \code{\textbackslash toprule}, \code{\textbackslash midrule}, and \code{\textbackslash bottomrule}; avoid a grid of vertical lines unless a journal/template specifically requires it.

\section{Table Notes with threeparttable}

\sourcehead
\begin{lstlisting}
\begin{table}[htbp]
\centering
\caption{Main Results}
\begin{threeparttable}
\begin{tabular}{lcc}
\toprule
& (1) & (2) \\
\midrule
Variable of interest & 0.125 & 0.110 \\
                     & (0.041) & (0.045) \\
\bottomrule
\end{tabular}
\begin{tablenotes}
\footnotesize
\item \textit{Notes:} Standard errors are in parentheses.
\end{tablenotes}
\end{threeparttable}
\end{table}
\end{lstlisting}

\resulthead
\begin{table}[H]
\centering
\caption{Main Results}
\begin{threeparttable}
\begin{tabular}{lcc}
\toprule
& (1) & (2) \\
\midrule
Variable of interest & 0.125 & 0.110 \\
                     & (0.041) & (0.045) \\
\bottomrule
\end{tabular}
\begin{tablenotes}
\footnotesize
\item \textit{Notes:} Standard errors are in parentheses.
\end{tablenotes}
\end{threeparttable}
\end{table}

\section{Panels and cmidrule}

\sourcehead
\begin{lstlisting}
\begin{tabular}{lcc}
\toprule
& (1) & (2) \\
\cmidrule(lr){2-3}
& Outcome A & Outcome B \\
\midrule
\multicolumn{3}{l}{\textit{Panel A: Full sample}} \\
Variable & 0.10 & 0.20 \\
\addlinespace
\multicolumn{3}{l}{\textit{Panel B: Subsample}} \\
Variable & 0.08 & 0.16 \\
\bottomrule
\end{tabular}
\end{lstlisting}

\resulthead
\begin{center}
\begin{tabular}{lcc}
\toprule
& (1) & (2) \\
\cmidrule(lr){2-3}
& Outcome A & Outcome B \\
\midrule
\multicolumn{3}{l}{\textit{Panel A: Full sample}} \\
Variable & 0.10 & 0.20 \\
\addlinespace
\multicolumn{3}{l}{\textit{Panel B: Subsample}} \\
Variable & 0.08 & 0.16 \\
\bottomrule
\end{tabular}
\end{center}

\section{What Does [htbp] Mean?}

For \code{\textbackslash begin\{table\}[htbp]} or \code{figure}, the letters suggest acceptable float positions:

\begin{itemize}
  \item \code{h}: approximately here;
  \item \code{t}: top of a page;
  \item \code{b}: bottom of a page;
  \item \code{p}: dedicated float page.
\end{itemize}

LaTeX is allowed to move tables and figures to improve page layout.

\warning{Do not use \code{[H]} on every table and figure just because a float moved. Use it selectively. Too many forced floats can create ugly white space.}

\section{Common Table Errors}

\subsection{Wrong number of ampersands}

If the table declares three columns, each row should normally contain two \code{\&} separators.

\begin{lstlisting}
\begin{tabular}{lcc}
Variable & Mean & SD \\       % correct: 3 columns
Income & 50.2 \\              % wrong: only 2 entries
\end{tabular}
\end{lstlisting}

The compiler may report \textbf{Extra alignment tab} or a confusing alignment error.

\subsection{Forgetting row endings}

\begin{lstlisting}
Income & 50.2 & 12.4 \\
Age    & 39.7 & 11.3 \\
\end{lstlisting}

Each completed row normally ends with \code{\textbackslash\textbackslash}.

\subsection{A table is too wide}

First, reduce unnecessary columns or shorten headings. If needed, use a smaller font locally:

\begin{lstlisting}
\begin{table}[htbp]
\centering
\small
...
\end{table}
\end{lstlisting}

Avoid shrinking important tables until the text becomes unreadable.

% ============================================================
\part{Figures and Graphics}
% ============================================================

\section{Your First Figure}

Load \code{graphicx} in the preamble. Then:

\sourcehead
\begin{lstlisting}
\begin{figure}[htbp]
\centering
\includegraphics[width=0.8\textwidth]{figures/main_result.pdf}
\caption{Main Result}
\label{fig:main}
\end{figure}
\end{lstlisting}

Because this handbook must compile without external files, the compiled example below uses a placeholder:

\resulthead
\begin{figure}[H]
\centering
\fbox{\parbox[c][3.2cm][c]{0.72\textwidth}{\centering Your exported R/Stata/Python figure would appear here.}}
\caption{Main Result}
\label{fig:placeholder-demo}
\end{figure}

\section{Choosing Figure File Types}

\begin{itemize}
  \item \textbf{PDF}: excellent for vector plots exported from statistical software.
  \item \textbf{PNG}: good for screenshots and raster graphics.
  \item \textbf{JPG}: best for photographs, not usually statistical plots.
\end{itemize}

For line graphs and coefficient plots, vector PDF usually stays sharp when zoomed.

\section{Sizing and Cropping Figures}

\sourcehead
\begin{lstlisting}
\includegraphics[width=0.75\textwidth]{figure.pdf}
\end{lstlisting}

Common alternatives:

\begin{lstlisting}
\includegraphics[width=\textwidth]{figure.pdf}
\includegraphics[height=7cm]{figure.pdf}
\includegraphics[trim=10 10 10 10,clip,width=0.8\textwidth]{figure.pdf}
\end{lstlisting}

The four trim numbers mean left, bottom, right, top (in TeX points unless units are supplied).

\section{Figure File Not Found}

If LaTeX says it cannot find the image, check:

\begin{enumerate}
  \item exact filename;
  \item exact extension;
  \item whether the file is in the expected folder;
  \item capitalization (important on some systems);
  \item whether the filename contains awkward characters.
\end{enumerate}

A safe project layout is:

\begin{lstlisting}
project/
  main.tex
  figures/
    fig1.pdf
    fig2.pdf
  tables/
  references.bib
\end{lstlisting}

\section{Keeping Floats from Crossing a Section}

If you want pending floats placed before the next section, load \code{placeins} and use:

\begin{lstlisting}
\FloatBarrier
\end{lstlisting}

Use it selectively rather than after every paragraph.

% ============================================================
\part{Cross-References, Footnotes, Links, and Citations}
% ============================================================

\section{Labels and References}

Give important objects a label and let LaTeX insert the current number.

\begin{center}
\begin{tabularx}{0.95\textwidth}{@{}lll@{}}
\toprule
Object & Label convention & Refer with \\
\midrule
Equation & \code{eq:baseline} & \code{Equation~\textbackslash eqref\{eq:baseline\}} \\
Table & \code{tab:main} & \code{Table~\textbackslash ref\{tab:main\}} \\
Figure & \code{fig:event} & \code{Figure~\textbackslash ref\{fig:event\}} \\
Section & \code{sec:data} & \code{Section~\textbackslash ref\{sec:data\}} \\
Appendix & \code{app:proofs} & \code{Appendix~\textbackslash ref\{app:proofs\}} \\
\bottomrule
\end{tabularx}
\end{center}

The prefixes are for humans, not LaTeX. They make labels easier to manage.

\section{Where Should the label Go?}

For figures and tables, put \code{\textbackslash label} \textbf{after} \code{\textbackslash caption}.

\begin{lstlisting}
\caption{Main Results}
\label{tab:main}
\end{lstlisting}

For equations, put the label inside the equation environment.

\warning{If a table reference shows the section number instead of the table number, your \code{\textbackslash label} may be before the caption.}

\section{Footnotes}

\sourcehead
\begin{lstlisting}
This restriction is standard.\footnote{A longer explanation can go here.}
\end{lstlisting}
\resulthead
This restriction is standard.\footnote{A longer explanation can go here.}

LaTeX numbers footnotes automatically.

\section{URLs and Hyperlinks}

With \code{hyperref} loaded:

\sourcehead
\begin{lstlisting}
\url{https://www.nber.org/}

\href{https://www.nber.org/}{NBER website}
\end{lstlisting}

\resulthead
\url{https://www.nber.org/}

\href{https://www.nber.org/}{NBER website}

Use \code{\textbackslash url\{...\}} for long URLs because it handles line breaking better than plain text.

\section{Citations: natbib Basics}

With \code{natbib}, the two commands you will use constantly are:

\begin{center}
\begin{tabularx}{0.94\textwidth}{@{}lXX@{}}
\toprule
Command & Use & Typical output \\
\midrule
\code{\textbackslash citet\{key\}} & Author is part of your sentence & Author (2020) \\
\code{\textbackslash citep\{key\}} & Entire citation is parenthetical & (Author, 2020) \\
\bottomrule
\end{tabularx}
\end{center}

\sourcehead
\begin{lstlisting}
\citet{sample2020} develops the framework.
A related literature studies this question \citep{sample2020}.
\end{lstlisting}

For this handbook, we will not execute those example keys, but in a real paper they refer to entries in \code{references.bib}.

\section{Your references.bib File}

A journal article entry might look like:

\begin{lstlisting}
@article{sample2020,
  author  = {Author, Alice and Scholar, Bob},
  title   = {A Sample Economics Paper},
  journal = {Journal of Economic Examples},
  year    = {2020},
  volume  = {10},
  number  = {2},
  pages   = {100--125},
  doi     = {10.xxxx/example}
}
\end{lstlisting}

At the end of your main \code{.tex} file:

\begin{lstlisting}
\bibliographystyle{plainnat}
\bibliography{references}
\end{lstlisting}

Do not write \code{references.bib} inside \code{\textbackslash bibliography\{\}}; write only the base name \code{references}.

\section{Citation Variants You May Need}

\begin{lstlisting}
\citep{key1,key2}        % multiple papers
\citep[see][]{key}       % (see Author, 2020)
\citep[p.~12]{key}       % page citation
\citet*{key}             % full author list (style-dependent)
\end{lstlisting}

\warning{If citations show as question marks, check the citation key spelling, bibliography commands, and compile cycle before rewriting everything.}

% ============================================================
\part{Economics Paper-Writing Patterns You Will Actually Reuse}
% ============================================================

\section{Front Matter: Title, Affiliation, Contact, and Thanks}

Economics papers often need more information than a simple name under the title. A common pattern is to put the affiliation directly under the author name and use \code{\textbackslash thanks} for contact information or acknowledgments.

\sourcehead
\begin{lstlisting}
\title{Why Does X Affect Y?}

\author{
Alice Researcher\thanks{Department of Economics,
Example University. Email: alice@example.edu.}
}

\date{September 2026}

\begin{document}
\maketitle
\end{document}
\end{lstlisting}

\resulthead
\begin{center}
{\Large\bfseries Why Does X Affect Y?}\par
\vspace{0.7em}
Alice Researcher\textsuperscript{*}\par
\vspace{0.4em}
September 2026
\end{center}
\smallskip
\noindent\textsuperscript{*}\footnotesize Department of Economics, Example University. Email: alice@example.edu.
\normalsize

\whyhead
\begin{itemize}
  \item \code{\textbackslash thanks\{...\}} creates a title-page footnote.
  \item You can use it for affiliation, email, acknowledgments, funding, or a disclosure statement.
  \item For a working paper, keep the title block simple unless your department or conference has a required format.
\end{itemize}

\warning{Do not manually type an asterisk and then manually create a footnote. Let \code{\textbackslash thanks} keep the marker and note connected.}

\section{Two or More Authors}

For an early draft, the simplest approach is often enough:

\sourcehead
\begin{lstlisting}
\author{
Alice Researcher\\Example University
\and
Bob Economist\\Another University
}
\end{lstlisting}

\resulthead
\begin{center}
Alice Researcher \hspace{3em} Bob Economist\\
Example University \hspace{3em} Another University
\end{center}

The command \code{\textbackslash and} tells the standard \code{article} class to separate authors. More elaborate journal title pages may require a journal template or a package, but beginners should first learn the standard mechanism.

\section{Abstract, Keywords, and JEL Codes}

A compact economics-paper front matter block often includes an abstract followed by keywords and JEL classifications.

\sourcehead
\begin{lstlisting}
\begin{abstract}
This paper studies how X affects Y using newly assembled data.
The analysis documents three main findings and discusses their
implications for Z.
\end{abstract}

\noindent\textbf{Keywords:} labor markets, firms, productivity

\noindent\textbf{JEL codes:} J24, J31, O33
\end{lstlisting}

\resulthead
\begin{abstract}
This paper studies how X affects Y using newly assembled data. The analysis documents three main findings and discusses their implications for Z.
\end{abstract}

\noindent\textbf{Keywords:} labor markets, firms, productivity

\noindent\textbf{JEL codes:} J24, J31, O33

\tip{The exact JEL codes belong to your research topic. The LaTeX lesson here is simply how to typeset a compact metadata block.}

\section{Writing the Introduction with LaTeX, Not Fighting LaTeX}

Your introduction is mostly ordinary prose. Do not put ordinary words into math mode, and do not add manual line breaks after every sentence.

\sourcehead
\begin{lstlisting}
\section{Introduction}
\label{sec:introduction}

How does X affect Y? This question matters because ...

This paper makes three contributions. First, ... Second, ...
Third, ...

The rest of the paper is organized as follows. Section~\ref{sec:data}
describes the data, Section~\ref{sec:design} explains the research
design, and Section~\ref{sec:results} presents the main results.
\end{lstlisting}

\resulthead
\textbf{1 Introduction}

How does X affect Y? This question matters because ...

This paper makes three contributions. First, ... Second, ... Third, ...

The rest of the paper is organized as follows. Section 3 describes the data, Section 4 explains the research design, and Section 5 presents the main results.

\tip{Use cross-references even while drafting. If sections move later, LaTeX updates the numbers automatically.}

\section{Defining Variables Clearly in Prose}

Economics papers repeatedly introduce outcomes, treatments, controls, prices, quantities, indices, and sample indicators. The cleanest habit is to define notation in a sentence immediately after it first appears.

\sourcehead
\begin{lstlisting}
Let $Y_{it}$ denote the outcome for unit $i$ in period $t$.
The variable $X_{it}$ is the main explanatory variable, and
$\mathbf{W}_{it}$ is a vector of observed controls.
The sample contains $N=12{,}480$ unit-period observations.
\end{lstlisting}

\resulthead
Let $Y_{it}$ denote the outcome for unit $i$ in period $t$.
The variable $X_{it}$ is the main explanatory variable, and
$\mathbf{W}_{it}$ is a vector of observed controls.
The sample contains $N=12{,}480$ unit-period observations.

\whyhead
\begin{itemize}
  \item Use math mode only for symbols: \code{\$Y\_\{it\}\$}, not for the surrounding sentence.
  \item \code{\textbackslash mathbf\{W\}} makes a bold mathematical symbol for a vector.
  \item \code{12\{,\}480} protects the comma as punctuation rather than mathematical spacing.
\end{itemize}

\section{A Small Notation Style Guide}

Consistency matters more than any single notation choice. A common pattern is:

\begin{center}
\begin{tabularx}{0.94\textwidth}{@{}lllX@{}}
\toprule
Meaning & Source & Output & Typical use \\
\midrule
Scalar & \code{\$x\$} & $x$ & one variable or parameter \\
Subscripted observation & \code{\$x\_i\$} & $x_i$ & observation $i$ \\
Population parameter & \code{\$\textbackslash beta\$} & $\beta$ & unknown parameter \\
Estimate & \code{\$\textbackslash hat\{\textbackslash beta\}\$} & $\hat{\beta}$ & estimated parameter \\
Vector & \code{\$\textbackslash mathbf\{x\}\$} & $\mathbf{x}$ & vector of variables \\
Expectation & \code{\$\textbackslash E[Y]\$} & $\E[Y]$ & expected value \\
Indicator & \code{\$\textbackslash mathbf\{1\}\textbackslash\{A\textbackslash\}\$} & $\mathbf{1}\{A\}$ & indicator for event $A$ \\
\bottomrule
\end{tabularx}
\end{center}

\tip{Pick one convention for vectors and estimates and use it everywhere. Inconsistent notation is much harder for a reader than simple notation.}

\section{Displaying a Baseline Specification Without Building a Formula Handbook}

You will often need one or two displayed equations even in an empirical paper. The important LaTeX skill is the presentation: number the equation, label it, and define its objects in prose.

\sourcehead
\begin{lstlisting}
\begin{equation}
Y_{it} = \alpha + \beta X_{it} + \mathbf{W}_{it}'\gamma + u_{it}.
\label{eq:baseline}
\end{equation}

In Equation~\eqref{eq:baseline}, $Y_{it}$ is the outcome,
$X_{it}$ is the variable of interest, and $\mathbf{W}_{it}$
contains the controls. The coefficient $\beta$ is the parameter
of interest.
\end{lstlisting}

\resulthead
\begin{equation}
Y_{it} = \alpha + \beta X_{it} + \mathbf{W}_{it}'\gamma + u_{it}.
\label{eq:paperpattern-baseline}
\end{equation}

In Equation~\eqref{eq:paperpattern-baseline}, $Y_{it}$ is the outcome,
$X_{it}$ is the variable of interest, and $\mathbf{W}_{it}$
contains the controls. The coefficient $\beta$ is the parameter of interest.

\tip{The purpose here is not to memorize a specific econometric model. It is to learn the paper-writing pattern: equation, label, reference, then definition.}

\section{Long Variable Names: Use Text or a Command}

Do not force long English phrases into math italics.

\sourcehead
\begin{lstlisting}
% Hard to read
$average_income_{county,t}$

% Better for a named variable
$\text{Income}_{ct}$

% Or define a reusable command in the preamble
\newcommand{\Income}{\mathrm{Income}}
...
$\Income_{ct}$
\end{lstlisting}

\resulthead
Better typesetting: $\text{Income}_{ct}$ or $\mathrm{Income}_{ct}$.

\whyhead
\code{\textbackslash text\{Income\}} uses normal text inside math. \code{\textbackslash mathrm\{Income\}} uses upright Roman letters and is useful for short variable names that behave like symbols.

\section{Assumptions, Definitions, Propositions, and Proofs}

Theory papers use these constantly, but empirical papers may also need formal assumptions or definitions. The \code{amsthm} package handles numbering and typography.

In the preamble:

\begin{lstlisting}
\usepackage{amsthm}

\newtheorem{assumption}{Assumption}
\newtheorem{proposition}{Proposition}
\theoremstyle{definition}
\newtheorem{definition}{Definition}
\end{lstlisting}

Then in the paper:

\sourcehead
\begin{lstlisting}
\begin{assumption}[Sampling]
Observations are independently sampled from the population.
\end{assumption}

\begin{proposition}
Under Assumption~1, the object of interest is identified.
\end{proposition}

\begin{proof}
The argument follows in three steps. First, ... Second, ...
Finally, ...
\end{proof}
\end{lstlisting}

\resulthead
\begin{assumption}[Sampling]
Observations are independently sampled from the population.
\end{assumption}

\begin{proposition}
Under Assumption~1, the object of interest is identified.
\end{proposition}

\begin{proof}
The argument follows in three steps. First, ... Second, ... Finally, ...
\end{proof}

\warning{Do not manually type ``Assumption 1'' or ``Proposition 2.'' Automatic numbering prevents errors when you add or delete results.}

\section{Numbered Conditions: (i), (ii), (iii)}

Economics writing often lists assumptions, mechanisms, sample restrictions, or contributions using Roman numerals.

\sourcehead
\begin{lstlisting}
The analysis imposes three sample restrictions:
\begin{enumerate}[label=(\roman*)]
  \item observations must have nonmissing outcomes;
  \item the main covariate must be observed; and
  \item the unit must appear for at least two periods.
\end{enumerate}
\end{lstlisting}

\resulthead
The analysis imposes three sample restrictions:
\begin{enumerate}[label=(\roman*)]
  \item observations must have nonmissing outcomes;
  \item the main covariate must be observed; and
  \item the unit must appear for at least two periods.
\end{enumerate}

The \code{enumitem} package makes labels such as \code{(i)}, \code{(a)}, or \code{1.} easy to control.

\section{Summary Statistics That Look Like an Economics Paper}

A useful summary-statistics table usually needs readable variable names, units when necessary, descriptive moments, and the number of observations.

\sourcehead
\begin{lstlisting}
\begin{table}[htbp]
\centering
\caption{Summary Statistics}
\label{tab:summary-econ}
\begin{threeparttable}
\begin{tabular}{lrrrrr}
\toprule
& Mean & SD & P25 & Median & N \\
\midrule
Outcome index       & 0.42 & 0.18 & 0.29 & 0.41 & 12,480 \\
Annual income (USD) & 54,210 & 18,430 & 41,500 & 51,900 & 12,480 \\
Age (years)         & 39.6 & 11.2 & 31 & 39 & 12,480 \\
\bottomrule
\end{tabular}
\begin{tablenotes}
\footnotesize
\item \textit{Notes:} The table reports descriptive statistics for
      the estimation sample. Monetary values are in 2026 dollars.
\end{tablenotes}
\end{threeparttable}
\end{table}
\end{lstlisting}

\resulthead
\begin{table}[H]
\centering
\caption{Summary Statistics}
\begin{threeparttable}
\begin{tabular}{lrrrrr}
\toprule
& Mean & SD & P25 & Median & N \\
\midrule
Outcome index       & 0.42 & 0.18 & 0.29 & 0.41 & 12,480 \\
Annual income (USD) & 54,210 & 18,430 & 41,500 & 51,900 & 12,480 \\
Age (years)         & 39.6 & 11.2 & 31 & 39 & 12,480 \\
\bottomrule
\end{tabular}
\begin{tablenotes}
\footnotesize
\item \textit{Notes:} The table reports descriptive statistics for the estimation sample. Monetary values are in 2026 dollars.
\end{tablenotes}
\end{threeparttable}
\end{table}

\tip{The note should tell the reader what sample the table uses and clarify non-obvious units or transformations.}

\section{A Regression-Style Results Table}

Even if Stata or R eventually generates your tables, it helps to understand the underlying LaTeX structure. Economics results tables commonly place estimates on one row and standard errors on the next, followed by specification information.

\sourcehead
\begin{lstlisting}
\begin{table}[htbp]
\centering
\caption{Main Results}
\label{tab:main-econ}
\begin{threeparttable}
\begin{tabular}{lccc}
\toprule
& (1) & (2) & (3) \\
& Outcome A & Outcome A & Outcome B \\
\midrule
Variable of interest & 0.118*** & 0.104*** & 0.072** \\
                     & (0.031)  & (0.034)  & (0.029) \\
\addlinespace
Controls             & No  & Yes & Yes \\
Unit fixed effects   & No  & Yes & Yes \\
Time fixed effects   & No  & Yes & Yes \\
Observations         & 12,480 & 12,480 & 12,480 \\
$R^2$                & 0.18 & 0.31 & 0.27 \\
\bottomrule
\end{tabular}
\begin{tablenotes}
\footnotesize
\item \textit{Notes:} Standard errors are in parentheses.
      * $p<0.10$, ** $p<0.05$, *** $p<0.01$.
\end{tablenotes}
\end{threeparttable}
\end{table}
\end{lstlisting}

\resulthead
\begin{table}[H]
\centering
\caption{Main Results}
\begin{threeparttable}
\begin{tabular}{lccc}
\toprule
& (1) & (2) & (3) \\
& Outcome A & Outcome A & Outcome B \\
\midrule
Variable of interest & 0.118*** & 0.104*** & 0.072** \\
                     & (0.031)  & (0.034)  & (0.029) \\
\addlinespace
Controls             & No  & Yes & Yes \\
Unit fixed effects   & No  & Yes & Yes \\
Time fixed effects   & No  & Yes & Yes \\
Observations         & 12,480 & 12,480 & 12,480 \\
$R^2$                & 0.18 & 0.31 & 0.27 \\
\bottomrule
\end{tabular}
\begin{tablenotes}
\footnotesize
\item \textit{Notes:} Standard errors are in parentheses. * $p<0.10$, ** $p<0.05$, *** $p<0.01$.
\end{tablenotes}
\end{threeparttable}
\end{table}

\whyhead
The first column is left aligned because it contains text. Result columns are centered here for simplicity. In a production table, packages such as \code{siunitx} can align numbers on decimal points.

\section{Significance Stars Without Retyping Them}

If you manually build a small table, a helper command can keep star formatting consistent.

In the preamble:

\begin{lstlisting}
\newcommand{\sym}[1]{\ifmmode^{#1}\else\(^{#1}\)\fi}
\end{lstlisting}

Then:

\sourcehead
\begin{lstlisting}
0.118\sym{***}
\end{lstlisting}

\resulthead
$0.118^{***}$

\tip{If a statistical package generates the table automatically, let that package handle stars rather than adding them manually afterward.}

\section{Aligning Numbers on the Decimal Point}

For tables with many numerical columns, \code{siunitx} can improve alignment.

In the preamble:

\begin{lstlisting}
\usepackage{siunitx}
\end{lstlisting}

\sourcehead
\begin{lstlisting}
\begin{tabular}{l S[table-format=1.3] S[table-format=2.2]}
\toprule
Variable & {Estimate} & {Mean} \\
\midrule
A & 0.125 & 4.20 \\
B & -0.043 & 12.75 \\
\bottomrule
\end{tabular}
\end{lstlisting}

\resulthead
\begin{center}
\begin{tabular}{l S[table-format=-1.3] S[table-format=2.2]}
\toprule
Variable & {Estimate} & {Mean} \\
\midrule
A & 0.125 & 4.20 \\
B & -0.043 & 12.75 \\
\bottomrule
\end{tabular}
\end{center}

The \code{S} column type understands numbers and aligns them at the decimal point.

\section{Column Groups with multicolumn and cmidrule}

Economics tables often compare outcomes or samples in grouped columns.

\sourcehead
\begin{lstlisting}
\begin{tabular}{lcccc}
\toprule
& \multicolumn{2}{c}{Full Sample}
& \multicolumn{2}{c}{Subsample} \\
\cmidrule(lr){2-3}\cmidrule(lr){4-5}
& (1) & (2) & (3) & (4) \\
\midrule
Variable & 0.10 & 0.12 & 0.08 & 0.09 \\
\bottomrule
\end{tabular}
\end{lstlisting}

\resulthead
\begin{center}
\begin{tabular}{lcccc}
\toprule
& \multicolumn{2}{c}{Full Sample}
& \multicolumn{2}{c}{Subsample} \\
\cmidrule(lr){2-3}\cmidrule(lr){4-5}
& (1) & (2) & (3) & (4) \\
\midrule
Variable & 0.10 & 0.12 & 0.08 & 0.09 \\
\bottomrule
\end{tabular}
\end{center}

\whyhead
\code{\textbackslash multicolumn\{2\}\{c\}\{Full Sample\}} makes one heading span two columns. \code{\textbackslash cmidrule(lr)\{2-3\}} draws a short rule under those two columns.

\section{Importing a Table Generated by Stata or R}

A productive research workflow often separates estimation from document formatting. Your statistical software writes a small \code{.tex} table body, and the paper imports it.

A project can look like:

\begin{lstlisting}
paper/
  main.tex
  references.bib
  tables/
    main_results.tex
    summary_stats.tex
  figures/
    event_plot.pdf
\end{lstlisting}

Then your paper can use:

\sourcehead
\begin{lstlisting}
\begin{table}[htbp]
\centering
\caption{Main Results}
\label{tab:main-imported}
\input{tables/main_results.tex}
\end{table}
\end{lstlisting}

\resulthead
The compiled PDF shows the contents of \code{tables/main\_results.tex} exactly where \code{\textbackslash input} appears, while \code{main.tex} stays short and readable.

\tip{Decide whether your exported file contains only the \code{tabular} material or an entire \code{table} environment. Do not accidentally wrap a complete \code{table} inside another \code{table}.}

\section{A Figure with a Professional Note}

Figures in economics papers often need a note explaining the sample, construction, units, or confidence intervals.

\sourcehead
\begin{lstlisting}
\begin{figure}[htbp]
\centering
\includegraphics[width=0.78\textwidth]{figures/main_figure.pdf}
\caption{Outcome by Group}
\label{fig:main-econ}

\begin{minipage}{0.88\textwidth}
\footnotesize
\textit{Notes:} The figure plots group means by year.
Vertical bars denote 95\% confidence intervals.
\end{minipage}
\end{figure}
\end{lstlisting}

\resulthead
\begin{figure}[H]
\centering
\fbox{\rule{0pt}{1.55in}\rule{0.72\textwidth}{0pt}}
\caption{Outcome by Group}
\begin{minipage}{0.88\textwidth}
\footnotesize
\textit{Notes:} The figure plots group means by year. Vertical bars denote 95\% confidence intervals.
\end{minipage}
\end{figure}

\tip{A caption says what the figure is. A note tells the reader what is necessary to interpret it correctly.}

\section{Two-Panel Figures with subcaption}

Multi-panel figures are common for heterogeneity, mechanisms, robustness, or different outcomes.

In the preamble:

\begin{lstlisting}
\usepackage{subcaption}
\end{lstlisting}

\sourcehead
\begin{lstlisting}
\begin{figure}[htbp]
\centering
\begin{subfigure}[t]{0.47\textwidth}
  \centering
  \includegraphics[width=\linewidth]{figures/panel_a.pdf}
  \caption{Outcome A}
  \label{fig:panel-a}
\end{subfigure}
\hfill
\begin{subfigure}[t]{0.47\textwidth}
  \centering
  \includegraphics[width=\linewidth]{figures/panel_b.pdf}
  \caption{Outcome B}
  \label{fig:panel-b}
\end{subfigure}
\caption{Effects Across Outcomes}
\label{fig:two-panel}
\end{figure}
\end{lstlisting}

\resulthead
\begin{figure}[H]
\centering
\begin{subfigure}[t]{0.45\textwidth}
  \centering
  \fbox{\rule{0pt}{1.0in}\rule{0.9\linewidth}{0pt}}
  \caption{Outcome A}
\end{subfigure}
\hfill
\begin{subfigure}[t]{0.45\textwidth}
  \centering
  \fbox{\rule{0pt}{1.0in}\rule{0.9\linewidth}{0pt}}
  \caption{Outcome B}
\end{subfigure}
\caption{Effects Across Outcomes}
\end{figure}

\whyhead
Each panel is a \code{subfigure}. The entire object is still one numbered figure with one overall caption.

\section{Referencing Panels, Tables, and Appendices in Sentences}

Use the tilde \code{\textasciitilde} before a reference so the label and number do not split across lines.

\sourcehead
\begin{lstlisting}
Table~\ref{tab:main} reports the baseline estimates.
Figure~\ref{fig:main} shows the corresponding pattern visually.
Appendix~\ref{app:robustness} reports additional robustness checks.
Equation~\eqref{eq:baseline} defines the main specification.
\end{lstlisting}

\resulthead
Table 2 reports the baseline estimates. Figure 3 shows the corresponding pattern visually. Appendix A reports additional robustness checks. Equation (1) defines the main specification.

\tip{Never type ``Table 2'' manually in the source if the number can change.}

\section{A Variable-Definition or Data-Construction Table}

For complicated data work, a compact variable table can be more readable than a long paragraph.

\sourcehead
\begin{lstlisting}
\begin{table}[htbp]
\centering
\caption{Key Variable Definitions}
\begin{tabularx}{0.95\textwidth}{@{}lXl@{}}
\toprule
Variable & Definition & Unit \\
\midrule
Outcome & Primary outcome constructed from administrative records. & Index \\
Income  & Annual pre-tax labor income. & 2026 USD \\
Age     & Age at the end of the calendar year. & Years \\
\bottomrule
\end{tabularx}
\end{table}
\end{lstlisting}

\resulthead
\begin{center}
\begin{tabularx}{0.92\textwidth}{@{}lXl@{}}
\toprule
Variable & Definition & Unit \\
\midrule
Outcome & Primary outcome constructed from administrative records. & Index \\
Income  & Annual pre-tax labor income. & 2026 USD \\
Age     & Age at the end of the calendar year. & Years \\
\bottomrule
\end{tabularx}
\end{center}

\section{Showing Sample Construction as a Small Table}

Sample construction is often clearer when the reader can see how many observations remain after each restriction.

\sourcehead
\begin{lstlisting}
\begin{tabular}{lr}
\toprule
Step & Observations \\
\midrule
Raw administrative records & 18,420 \\
Drop missing outcome        & 16,905 \\
Apply age restriction       & 14,212 \\
Require complete covariates & 12,480 \\
\bottomrule
\end{tabular}
\end{lstlisting}

\resulthead
\begin{center}
\begin{tabular}{lr}
\toprule
Step & Observations \\
\midrule
Raw administrative records & 18,420 \\
Drop missing outcome        & 16,905 \\
Apply age restriction       & 14,212 \\
Require complete covariates & 12,480 \\
\bottomrule
\end{tabular}
\end{center}

This is a good example of using LaTeX to improve exposition rather than adding more mathematics.

\section{Robustness and Heterogeneity: Organize the Presentation}

A robustness section often becomes hard to read because it contains many specifications. LaTeX helps if you organize the prose and references consistently.

\sourcehead
\begin{lstlisting}
\section{Robustness and Heterogeneity}
\label{sec:robustness}

\subsection{Alternative Samples}
Table~\ref{tab:robust-samples} repeats the analysis using ...

\subsection{Alternative Variable Definitions}
Table~\ref{tab:robust-definitions} shows that ...

\subsection{Heterogeneity}
Figure~\ref{fig:heterogeneity} reports estimates by ...
\end{lstlisting}

\resulthead
\textbf{6 Robustness and Heterogeneity}

\textbf{6.1 Alternative Samples}\quad Table 6 repeats the analysis using ...

\textbf{6.2 Alternative Variable Definitions}\quad Table 7 shows that ...

\textbf{6.3 Heterogeneity}\quad Figure 4 reports estimates by ...

\tip{The LaTeX skill is simple sectioning and cross-referencing. The research judgment is deciding what belongs in the main text versus the appendix.}

\section{Appendix Table and Figure Numbering: A1, A2, ...}

In the standard \code{article} class, \code{\textbackslash appendix} changes section numbers to letters, but table and figure counters may continue from the main text. Many economics papers instead use labels such as Table A1 and Figure A1.

\sourcehead
\begin{lstlisting}
\clearpage
\appendix

\setcounter{table}{0}
\renewcommand{\thetable}{A\arabic{table}}
\setcounter{figure}{0}
\renewcommand{\thefigure}{A\arabic{figure}}

\section{Additional Results}

\begin{table}[htbp]
\centering
\caption{Additional Results}
\label{tab:appendix-main}
...
\end{table}
\end{lstlisting}

\resulthead
The first appendix table is numbered \textbf{Table A1}; the next is \textbf{Table A2}. The first appendix figure is \textbf{Figure A1}.

\warning{If your journal class already controls appendix numbering, follow the journal template instead of overriding it manually.}

\section{Appendix Numbering by Section: A.1, A.2, B.1, ...}

For a long appendix, you may want table numbers to include the appendix section letter.

\begin{lstlisting}
\appendix
\numberwithin{table}{section}
\numberwithin{figure}{section}
\end{lstlisting}

With appendix sections A and B, this can produce numbers such as Table A.1 and Table B.1.

\tip{Choose one appendix numbering convention and keep it consistent across tables, figures, and references.}

\section{Wide Tables: First Fix the Content, Then the Layout}

A wide table is common in empirical work. Do not immediately shrink it to unreadable text.

Try, in this order:

\begin{enumerate}
  \item shorten column headings;
  \item move details into the table note;
  \item remove redundant columns;
  \item use \code{\textbackslash small} locally;
  \item use a landscape page only if the table genuinely needs the width.
\end{enumerate}

For a landscape page, add:

\begin{lstlisting}
\usepackage{pdflscape}
\end{lstlisting}

Then:

\begin{lstlisting}
\begin{landscape}
\begin{table}[p]
\centering
\caption{Wide Robustness Table}
...
\end{table}
\end{landscape}
\end{lstlisting}

The \code{pdflscape} package rotates the page for PDF viewing as well as printing.

\section{Do Not Use resizebox as the First Solution}

This compiles:

\begin{lstlisting}
\resizebox{\textwidth}{!}{%
  \begin{tabular}{lcccccccc}
  ...
  \end{tabular}
}
\end{lstlisting}

But a table can become too small to read. Use \code{\textbackslash resizebox} only after improving the table design itself.

\section{Common Economics Typesetting Details}

Small details appear constantly in drafts.

\begin{center}
\begin{tabularx}{0.95\textwidth}{@{}lXl@{}}
\toprule
You want & Source & Output \\
\midrule
five percent & \code{5\textbackslash\%} & 5\% \\
range of years & \code{2000--2010} & 2000--2010 \\
negative number & \code{\$-0.25\$} & $-0.25$ \\
p-value & \code{\$p<0.05\$} & $p<0.05$ \\
R-squared & \code{\$R\^2\$} & $R^2$ \\
sample size & \code{\$N=1\{,\}250\$} & $N=1{,}250$ \\
dollars & \code{\textbackslash\$500} & \$500 \\
95 percent CI & \code{95\textbackslash\% CI} & 95\% CI \\
Table reference & \code{Table\textasciitilde\textbackslash ref\{tab:main\}} & Table 2 \\
\bottomrule
\end{tabularx}
\end{center}

\warning{A literal \code{\$} starts or ends math mode. To print the currency symbol, type \code{\textbackslash\$}.}

\section{Percentage Points vs. Percent Signs}

This is mostly a writing issue, but LaTeX makes the typography easy:

\sourcehead
\begin{lstlisting}
The outcome rises by 5\%.
The treatment increases the participation rate by 5 percentage points.
\end{lstlisting}

\resulthead
The outcome rises by 5\%.
The treatment increases the participation rate by 5 percentage points.

Do not use the percent symbol when you mean percentage points simply because it is shorter to type.

\section{Citing Papers in Different Sentence Positions}

Citation commands should fit the grammar of the sentence.

\sourcehead
\begin{lstlisting}
\citet{author2020} studies a related mechanism.

A large literature studies this question \citep{author2020,author2022}.

The measurement follows \citet[p.~18]{author2020}.
\end{lstlisting}

\resulthead
Author (2020) studies a related mechanism.

A large literature studies this question (Author, 2020; Author, 2022).

The measurement follows Author (2020, p. 18).

\tip{Use \code{\textbackslash citet} when the author is grammatically part of your sentence; use \code{\textbackslash citep} when the citation itself is parenthetical.}

\section{Avoid Hard-Coding Names of Sections, Tables, and Equations}

Bad drafting habit:

\begin{lstlisting}
As shown in Table 4, ...
See Section 6 for robustness.
\end{lstlisting}

Better:

\begin{lstlisting}
As shown in Table~\ref{tab:main}, ...
See Section~\ref{sec:robustness} for robustness.
\end{lstlisting}

If Table 4 becomes Table 5 later, your sentence remains correct automatically.

\section{Useful Cross-Reference Naming Conventions}

A predictable label system saves time in a long project.

\begin{center}
\begin{tabularx}{0.92\textwidth}{@{}llX@{}}
\toprule
Object & Example label & Meaning \\
\midrule
Section & \code{sec:data} & data section \\
Equation & \code{eq:baseline} & baseline equation \\
Table & \code{tab:main} & main results table \\
Figure & \code{fig:heterogeneity} & heterogeneity figure \\
Appendix & \code{app:data} & data appendix \\
Proposition & \code{prop:main} & main proposition \\
Assumption & \code{assump:sampling} & sampling assumption \\
\bottomrule
\end{tabularx}
\end{center}

The prefix is only for you; LaTeX does not require it.

\section{Reusable Commands for Long Names}

If a phrase or notation appears repeatedly, define it once.

\sourcehead
\begin{lstlisting}
% Preamble
\newcommand{\mainoutcome}{employment rate}
\newcommand{\Y}{Y_{it}}

% Body
The main outcome is the \mainoutcome.
We denote it by $\Y$.
\end{lstlisting}

\resulthead
The main outcome is the employment rate. We denote it by $Y_{it}$.

\tip{Custom commands are most valuable when a notation or term may change later. One edit in the preamble can then update the entire paper.}

\section{Draft Notes That Cannot Accidentally Look Final}

During a long paper project, visible draft markers are useful.

In the preamble:

\begin{lstlisting}
\newcommand{\todo}[1]{\textcolor{red}{[TODO: #1]}}
\end{lstlisting}

Then:

\sourcehead
\begin{lstlisting}
The estimate remains stable across specifications.
\todo{Add appendix-table reference here.}
\end{lstlisting}

\resulthead
The estimate remains stable across specifications. \textcolor{red}{[TODO: Add appendix-table reference here.]}

\warning{Search the source and final PDF for \code{TODO} before circulating or submitting the paper.}

\section{Anonymous and Named Drafts from One File}

For conferences or journals that require anonymous submission, a simple switch can avoid maintaining two separate papers.

\begin{lstlisting}
% In the preamble
\newif\ifanonymous
\anonymoustrue      % switch to \anonymousfalse for named draft

\ifanonymous
  \author{Anonymous}
\else
  \author{Alice Researcher\\Example University}
\fi
\end{lstlisting}

This is more advanced than ordinary beginner syntax, but the idea is useful: keep one source of truth and switch only the information that must differ.

\section{Double Spacing and Line Numbers for Submission}

Some journals, workshops, or advisors request a referee-friendly draft.

In the preamble:

\begin{lstlisting}
\usepackage{setspace}
\usepackage{lineno}
\end{lstlisting}

In the document:

\begin{lstlisting}
\doublespacing
\linenumbers
\end{lstlisting}

For your normal working draft, you may prefer single spacing and no line numbers. Follow the actual submission instructions rather than assuming one format is universal.

\section{A Reproducible Folder Structure for an Economics Project}

A clean project structure makes LaTeX easier to maintain and makes replication easier later.

\begin{lstlisting}
project/
  paper/
    main.tex
    references.bib
    sections/
      introduction.tex
      data.tex
      design.tex
      results.tex
    tables/
      summary_stats.tex
      main_results.tex
    figures/
      main_figure.pdf
  code/
    01_clean_data.do
    02_analysis.do
  data/
    raw/
    derived/
\end{lstlisting}

Then \code{main.tex} can contain:

\begin{lstlisting}
\input{sections/introduction}
\input{sections/data}
\input{sections/design}
\input{sections/results}
\end{lstlisting}

\tip{LaTeX does not make a project reproducible by itself, but a predictable folder structure makes it much easier to connect code, tables, figures, and the written paper.}

\section{A More Realistic Economics-Paper Skeleton}

The following is not a rule for every paper. It is a useful beginner scaffold showing how the pieces in this handbook fit together.

\begin{lstlisting}
\documentclass[11pt]{article}

\usepackage[margin=1in]{geometry}
\usepackage{amsmath,amssymb,amsthm}
\usepackage{graphicx}
\usepackage{booktabs,threeparttable,tabularx}
\usepackage{subcaption}
\usepackage[round,authoryear]{natbib}
\usepackage[hidelinks]{hyperref}

\title{Paper Title}
\author{Your Name\\Department of Economics, Your University}
\date{\today}

\begin{document}
\maketitle

\begin{abstract}
Question, data or approach, main finding, and contribution.
\end{abstract}

\noindent\textbf{Keywords:} keyword 1, keyword 2, keyword 3

\noindent\textbf{JEL codes:} A00, B00

\section{Introduction}
\label{sec:introduction}

\section{Background or Institutional Setting}
\label{sec:background}

\section{Data}
\label{sec:data}

\section{Research Design}
\label{sec:design}

\begin{equation}
Y_{it}=\alpha+\beta X_{it}+\mathbf{W}_{it}'\gamma+u_{it}.
\label{eq:baseline-paper}
\end{equation}

Define every symbol in Equation~\eqref{eq:baseline-paper}.

\section{Results}
\label{sec:results}

Table~\ref{tab:main-paper} reports the main results.
Figure~\ref{fig:main-paper} provides a graphical summary.

\section{Robustness and Extensions}
\label{sec:robustness-paper}

\section{Conclusion}

\clearpage
\bibliographystyle{plainnat}
\bibliography{references}

\clearpage
\appendix
\section{Additional Results}
\label{app:additional}

\end{document}
\end{lstlisting}

\whyhead
This skeleton intentionally contains only one generic equation. The emphasis is on the writing system around the research: sections, labels, references, tables, figures, bibliography, and appendix.



% ============================================================
\part{Appendices, Project Organization, and Reuse}
% ============================================================

\section{Appendices}

\sourcehead
\begin{lstlisting}
\clearpage
\appendix

\section{Additional Results}
\label{app:results}

\section{Additional Figures}
\end{lstlisting}

After \code{\textbackslash appendix}, standard article sections are normally labeled with letters.

\section{Page Breaks: newpage vs. clearpage}

\begin{itemize}
  \item \code{\textbackslash newpage}: starts a new page.
  \item \code{\textbackslash clearpage}: places pending figures/tables first, then starts a new page.
\end{itemize}

For starting an appendix or references section after many floats, \code{\textbackslash clearpage} is often useful.

\section{Splitting a Large Paper into Files}

A dissertation chapter or long paper becomes easier to manage if sections are separated.

\begin{lstlisting}
project/
  main.tex
  sections/
    introduction.tex
    data.tex
    methods.tex
    results.tex
  figures/
  tables/
  references.bib
\end{lstlisting}

In \code{main.tex}:

\begin{lstlisting}
\begin{document}
\maketitle

\input{sections/introduction}
\input{sections/data}
\input{sections/methods}
\input{sections/results}

\end{document}
\end{lstlisting}

The section files usually contain only content, not another \code{\textbackslash documentclass} or \code{\textbackslash begin\{document\}}.

\section{Custom Commands: Avoid Re-Typing Long Notation}

Suppose you repeatedly write expectations and probabilities. Put in the preamble:

\begin{lstlisting}
\newcommand{\E}{\mathbb{E}}
\newcommand{\Prob}{\mathbb{P}}
\end{lstlisting}

Then:

\sourcehead
\begin{lstlisting}
$\E[Y\mid X]$ and $\Prob(A)$
\end{lstlisting}
\resulthead
$\E[Y\mid X]$ and $\Prob(A)$.

\tip{Custom commands are valuable when notation appears many times. Do not create a macro for every one-off symbol.}

\section{Draft Comments and Temporary Notes}

The simplest private note is a source comment:

\begin{lstlisting}
% TODO: explain identification more clearly.
\end{lstlisting}

You can also define a visible draft command:

\begin{lstlisting}
\newcommand{\todo}[1]{\textbf{[TODO: #1]}}
\end{lstlisting}

Then \code{\textbackslash todo\{rewrite this paragraph\}} prints a visible reminder. Remove or disable such notes before submission.

\section{Safe File Names}

Prefer:

\begin{lstlisting}
main.tex
references.bib
fig_event_study.pdf
tab_summary.tex
\end{lstlisting}

Avoid complicated punctuation and ambiguous file names. Spaces may work in modern systems, but simple names reduce cross-platform problems.

% ============================================================
\part{Common Mistakes and Debugging}
% ============================================================

\section{The Beginner Error Map}

\begin{center}
\begin{tabularx}{0.98\textwidth}{@{}p{0.28\textwidth}X@{}}
\toprule
Message or symptom & First thing to check \\
\midrule
\texttt{Undefined control sequence} & Typo in a command, or missing package. \\
\texttt{Missing \$ inserted} & Math syntax used outside math mode; unmatched dollar sign. \\
\texttt{Extra alignment tab \&} & Too many ampersands in a table/alignment row. \\
\texttt{Missing \} inserted} & Missing closing brace nearby. \\
\texttt{Runaway argument?} & Often a missing closing brace or environment end. \\
\texttt{Environment ... undefined} & Missing package or misspelled environment. \\
\texttt{... ended by \textbackslash end\{...\}} & Begin/end environment names do not match. \\
Reference displays \texttt{??} & Compile again; verify the label exists and spelling matches. \\
Citation displays \texttt{?} & Verify citation key and bibliography workflow. \\
\texttt{File ... not found} & Wrong path, filename, extension, or capitalization. \\
\texttt{Overfull \textbackslash hbox} & Something is too wide for the line/page. \\
Figure/table moved & Normal float behavior; placement is a request, not an absolute command. \\
\bottomrule
\end{tabularx}
\end{center}

\section{Mistake: Missing or Extra Braces}

Wrong:

\begin{lstlisting}
\frac{a}{b
\end{lstlisting}

Correct:

\begin{lstlisting}
\frac{a}{b}
\end{lstlisting}

When the error message looks strange, inspect the previous few lines. A missing brace can confuse LaTeX far beyond the original typo.

\section{Mistake: Math Command Outside Math Mode}

Wrong:

\begin{lstlisting}
The coefficient is \beta.
\end{lstlisting}

Correct:

\begin{lstlisting}
The coefficient is $\beta$.
\end{lstlisting}

\section{Mistake: Unescaped Percent, Ampersand, or Underscore}

Wrong prose:

\begin{lstlisting}
The response rate is 80%.
R&D spending increases.
data_clean.csv
\end{lstlisting}

Correct prose:

\begin{lstlisting}
The response rate is 80\%.
R\&D spending increases.
data\_clean.csv
\end{lstlisting}

\section{Mistake: Mismatched Environments}

Wrong:

\begin{lstlisting}
\begin{align}
a &= b+c
\end{equation}
\end{lstlisting}

Correct:

\begin{lstlisting}
\begin{align}
a &= b+c
\end{align}
\end{lstlisting}

\section{Mistake: Using Double Dollars for Display Math}

You may see old source files that use:

\begin{lstlisting}
$$ y = a + bx $$
\end{lstlisting}

For modern LaTeX documents, prefer:

\begin{lstlisting}
\[
y = a + bx
\]
\end{lstlisting}

or a numbered \code{equation} environment.

\section{Mistake: Using Line-Break Commands for Paragraph Spacing}

Wrong habit:

\begin{lstlisting}
First paragraph.\\
\\
Second paragraph.
\end{lstlisting}

Better:

\begin{lstlisting}
First paragraph.

Second paragraph.
\end{lstlisting}

Use \code{\textbackslash\textbackslash} mainly where a line break is structurally required, such as rows in tables or aligned equations.

\section{Mistake: Manually Typing Numbers}

Avoid:

\begin{lstlisting}
As shown in Table 3 ...
Equation (5) implies ...
\end{lstlisting}

Prefer labels:

\begin{lstlisting}
As shown in Table~\ref{tab:main} ...
Equation~\eqref{eq:model} implies ...
\end{lstlisting}

If you later insert a table earlier in the paper, the numbers update automatically.

\section{Mistake: caption and label in the Wrong Order}

Better:

\begin{lstlisting}
\caption{Main Results}
\label{tab:main}
\end{lstlisting}

For tables and figures, \code{\textbackslash label} normally comes after \code{\textbackslash caption}.

\section{Mistake: Forcing Every Float with [H]}

\code{[H]} can be useful, but it says ``put this float exactly here.'' If you force every table and figure, LaTeX may leave large blank areas.

Start with \code{[htbp]}. Use \code{[H]} only when you have a specific reason.

\section{Mistake: Very Long URLs or Unbreakable Text}

Use:

\begin{lstlisting}
\url{https://example.com/a/very/long/path}
\end{lstlisting}

instead of pasting a long raw URL into prose. For code/file paths, \code{\textbackslash path\{...\}} from \code{url/hyperref} can also help.

\section{Overfull and Underfull Boxes}

An \textbf{Overfull \textbackslash hbox} warning means some line extends beyond its intended width. Common causes:

\begin{itemize}
  \item a long URL;
  \item a wide equation;
  \item a table that is too wide;
  \item a long unbreakable filename or code string.
\end{itemize}

An \textbf{Underfull \textbackslash hbox} warning means LaTeX had difficulty spacing a line attractively. It is often cosmetic, but you should visually inspect the page.

\section{A Systematic Debugging Method}

When something breaks:

\begin{enumerate}
  \item Read the \textbf{first} meaningful error, not the 20 errors that follow it.
  \item Inspect the reported line and several lines above it.
  \item Check braces, dollar signs, \code{\textbackslash begin}/\code{\textbackslash end}, and special characters.
  \item Ask what you changed most recently.
  \item Comment out the suspicious block with \code{\%} and compile again.
  \item Restore pieces gradually until the error returns.
\end{enumerate}

This ``reduce the problem'' strategy is much faster than randomly changing commands.

% ============================================================
\part{Complete Beginner Templates}
% ============================================================

\section{Minimal Problem-Set Template}

\begin{lstlisting}
\documentclass[11pt]{article}

\usepackage[margin=1in]{geometry}
\usepackage{amsmath,amssymb}

\title{ECON 000: Problem Set 1}
\author{Your Name}
\date{\today}

\begin{document}
\maketitle

\section*{Question 1}

Write the explanation in prose first. Short notation can stay inline,
such as $\beta$ or $x_i$.

For an important displayed expression:
\[
y_i = \alpha + \beta x_i + \varepsilon_i.
\]

For multiple steps:
\begin{align*}
A &= B+C,\\
  &= D+E.
\end{align*}

\section*{Question 2}

Continue here.

\end{document}
\end{lstlisting}

\section{Minimal Economics Paper Template}

\begin{lstlisting}
\documentclass[11pt]{article}

\usepackage[margin=1in]{geometry}
\usepackage{amsmath,amssymb}
\usepackage{graphicx}
\usepackage{subcaption}
\usepackage{pdflscape}
\usepackage{booktabs,threeparttable}
\usepackage[round,authoryear]{natbib}
\usepackage[hidelinks]{hyperref}

\title{Paper Title}
\author{Your Name\\Department of Economics, Your University}
\date{\today}

\begin{document}
\maketitle

\begin{abstract}
One paragraph: question, approach, main finding, contribution.
\end{abstract}

\section{Introduction}
\label{sec:intro}

State the question, why it matters, what you do, what you find,
and how the paper contributes.

\section{Background}

\section{Data}
\label{sec:data}

\section{Research Design}

\begin{equation}
y_i = \alpha + \beta x_i + \varepsilon_i.
\label{eq:baseline}
\end{equation}

Define every object in Equation~\eqref{eq:baseline} in the text.

\section{Results}

Table~\ref{tab:main} reports the main result.

\begin{table}[htbp]
\centering
\caption{Main Results}
\label{tab:main}
\begin{threeparttable}
\begin{tabular}{lcc}
\toprule
& (1) & (2) \\
\midrule
Variable of interest & 0.125 & 0.110 \\
                     & (0.041) & (0.045) \\
\midrule
Observations & 1,000 & 1,000 \\
\bottomrule
\end{tabular}
\begin{tablenotes}
\footnotesize
\item \textit{Notes:} Explain the sample, standard errors,
and any notation needed to understand the table.
\end{tablenotes}
\end{threeparttable}
\end{table}

\section{Conclusion}

\clearpage
\appendix
\section{Additional Results}

\bibliographystyle{plainnat}
\bibliography{references}

\end{document}
\end{lstlisting}

\section{Multi-File Paper Template}

\textbf{main.tex}

\begin{lstlisting}
\documentclass[11pt]{article}
\usepackage[margin=1in]{geometry}
\usepackage{amsmath,amssymb}
\usepackage{graphicx,booktabs}
\usepackage[round,authoryear]{natbib}
\usepackage[hidelinks]{hyperref}

\title{Paper Title}
\author{Your Name}
\date{}

\begin{document}
\maketitle

\input{sections/abstract}
\input{sections/introduction}
\input{sections/data}
\input{sections/design}
\input{sections/results}
\input{sections/conclusion}

\clearpage
\bibliographystyle{plainnat}
\bibliography{references}
\end{document}
\end{lstlisting}

\textbf{sections/introduction.tex}

\begin{lstlisting}
\section{Introduction}
\label{sec:intro}

Write the introduction here.
\end{lstlisting}

% ============================================================
\part{Quick Reference for Daily Use}
% ============================================================

\section{The Commands to Learn First}

\begin{center}
\small
\begin{tabularx}{0.98\textwidth}{@{}lXl@{}}
\toprule
Command & What it does & Example result \\
\midrule
\code{\textbackslash section\{...\}} & section heading & numbered heading \\
\code{\textbackslash textbf\{...\}} & bold text & \textbf{bold} \\
\code{\$...\$} & inline math & $x_i$ \\
\code{\textbackslash[...\textbackslash]} & display math & centered equation \\
\code{\_} & subscript & $x_i$ \\
\code{\^} & superscript & $x^2$ \\
\code{\textbackslash frac\{a\}\{b\}} & fraction & $\frac{a}{b}$ \\
\code{\textbackslash sum} & summation & $\sum_i x_i$ \\
\code{\textbackslash begin\{align*\}} & aligned derivation & aligned lines \\
\code{\textbackslash text\{...\}} & words inside math & $x=1\ \text{if treated}$ \\
\code{\textbackslash begin\{table\}} & table float & numbered table \\
\code{\textbackslash begin\{tabular\}} & rows/columns & table body \\
\code{\textbackslash includegraphics} & image import & figure \\
\code{\textbackslash caption} & table/figure title & caption \\
\code{\textbackslash label} & target for reference & internal label \\
\code{\textbackslash ref} & table/figure/section reference & number \\
\code{\textbackslash eqref} & equation reference & (number) \\
\code{\textbackslash footnote} & footnote & footnote marker \\
\code{\textbackslash citet} & textual citation & Author (year) \\
\code{\textbackslash citep} & parenthetical citation & (Author, year) \\
\code{\textbackslash multicolumn} & grouped table heading & spans columns \\
\code{\textbackslash input} & import another .tex file & inserts file contents \\
\code{\textbackslash appendix} & begin appendix & lettered appendix sections \\
\code{\textbackslash newcommand} & reusable notation/text & custom command \\
\bottomrule
\end{tabularx}
\end{center}


\section{Spacing, Alignment, and Tiny Commands - Fast Lookup}

These are the small commands you will use constantly while writing. For most of them, you do not need to memorize the exact spacing size; remember the rough order and copy the command when needed.

\subsection*{Math spacing ladder}

\begin{center}
\small
\begin{tabularx}{0.98\textwidth}{@{}lXl@{}}
\toprule
Command & What it does & Compiled example \\
\midrule
\code{\textbackslash,} & tiny positive math space & $a\,b$ \\
\code{\textbackslash:} & small/medium math space & $a\:b$ \\
\code{\textbackslash;} & medium-large math space & $a\;b$ \\
\code{\textbackslash!} & small \emph{negative} math space & $a\!b$ \\
\code{\textbackslash quad} & large fixed space (about 1 em) & $a\quad b$ \\
\code{\textbackslash qquad} & extra-large fixed space (about 2 em) & $a\qquad b$ \\
\bottomrule
\end{tabularx}
\end{center}

A useful visual order is
\[
a\!b
\quad < \quad
a\,b
\quad < \quad
a\:b
\quad < \quad
a\;b
\quad < \quad
a\quad b
\quad < \quad
a\qquad b.
\]

A common economics-paper pattern is to separate an equation from an index or condition:

\begin{lstlisting}
Y_i = \alpha + \beta X_i + \varepsilon_i,
\qquad i=1,\ldots,n.
\end{lstlisting}

which compiles as
\[
Y_i = \alpha + \beta X_i + \varepsilon_i,
\qquad i=1,\ldots,n.
\]

\tip{Use \code{\textbackslash quad} or \code{\textbackslash qquad} for visual separation, not to fake equation alignment. If several equations should line up at their equals signs, use \code{align} with \code{\&}.}

\subsection*{Alignment and layout shortcuts}

\begin{center}
\small
\begin{tabularx}{0.98\textwidth}{@{}lXl@{}}
\toprule
Command & Use & Typical pattern \\
\midrule
\code{\&} & alignment point in \code{align} or column separator in tables & \code{Y\_i \&= a+bX\_i} \\
\code{\textbackslash\textbackslash} & end a row/line in \code{align}, \code{tabular}, etc. & \code{... \textbackslash\textbackslash} \\
\code{\textbackslash notag} & suppress one equation number in \code{align} & \code{... \textbackslash notag} \\
\code{\textasciitilde} & nonbreaking space & \code{Table\string~\textbackslash ref\{tab:main\}} \\
\code{\textbackslash hfill} & flexible horizontal space; pushes material apart & \code{Left \textbackslash hfill Right} \\
\code{\textbackslash hspace\{...\}} & exact horizontal space & \code{\textbackslash hspace\{1em\}} \\
\code{\textbackslash smallskip} & small vertical gap & between nearby blocks \\
\code{\textbackslash medskip} & medium vertical gap & between nearby blocks \\
\code{\textbackslash bigskip} & larger vertical gap & between nearby blocks \\
\code{\textbackslash vspace\{...\}} & exact vertical space & \code{\textbackslash vspace\{0.5em\}} \\
\code{\textbackslash noindent} & suppress paragraph indentation once & \code{\textbackslash noindent Text...} \\
\code{\textbackslash centering} & center material inside a figure/table/group & before \code{\textbackslash includegraphics} \\
\code{\textbackslash phantom\{...\}} & reserve invisible width/height & manual alignment trick \\
\bottomrule
\end{tabularx}
\end{center}

For aligned equations, remember this pattern:

\begin{lstlisting}
\begin{align}
Y_i &= \alpha + \beta X_i + \varepsilon_i, \\
E[\varepsilon_i\mid X_i] &= 0. \notag
\end{align}
\end{lstlisting}

The \code{\&} marks the alignment point; placing it immediately before \code{=} lines up the equals signs.

\subsection*{Tiny math-writing helpers worth remembering}

\begin{center}
\small
\begin{tabularx}{0.98\textwidth}{@{}lXl@{}}
\toprule
Command & Why it is useful & Example \\
\midrule
\code{\textbackslash text\{...\}} & normal words inside math & $D_i=1\ \text{if treated}$ \\
\code{\textbackslash mathrm\{...\}} & upright math letters/units & $10\,\mathrm{km}$ \\
\code{\textbackslash operatorname\{Var\}} & custom operator with proper spacing & $\operatorname{Var}(X)$ \\
\code{\textbackslash left(...\textbackslash right)} & automatic delimiter sizing & $\left(\frac{a}{b}\right)$ \\
\code{\textbackslash dfrac\{a\}\{b\}} & display-style fraction in inline math & $\dfrac{a}{b}$ \\
\code{\textbackslash lvert x\textbackslash rvert} & properly sized absolute value & $\lvert x\rvert$ \\
\code{\textbackslash lVert x\textbackslash rVert} & norm & $\lVert x\rVert$ \\
\code{\textbackslash underbrace} & annotate part of an expression & $\underbrace{x_1+x_2}_{\text{two terms}}$ \\
\code{\textbackslash substack} & stack conditions under sums/limits & $\sum_{\substack{i=1\\i\ne j}}^n x_i$ \\
\code{\textbackslash ldots} / \code{\textbackslash cdots} & ellipses in lists / centered math & $x_1,\ldots,x_n$; $A_1\cdots A_n$ \\
\bottomrule
\end{tabularx}
\end{center}

\subsection*{Show a command literally in your PDF}

If you want the PDF to display a command rather than execute it, use inline verbatim:

\begin{lstlisting}
Use \verb|\qquad| for a large horizontal space.
\end{lstlisting}

This prints the command \verb|\qquad| literally. With the \code{listings} package, \code{\textbackslash lstinline|...|} is another convenient option.

\subsection*{Fast ``I need to...'' lookup}

\begin{center}
\small
\begin{tabularx}{0.98\textwidth}{@{}Xl@{}}
\toprule
I need to... & Try this first \\
\midrule
add a tiny math space & \code{\textbackslash,} \\
add a medium math space & \code{\textbackslash;} \\
add a large math space & \code{\textbackslash quad} \\
add an even larger math space & \code{\textbackslash qquad} \\
slightly tighten math spacing & \code{\textbackslash!} \\
keep ``Table 3'' on one line & \code{Table\string~\textbackslash ref\{...\}} \\
push one item to the far right & \code{\textbackslash hfill} \\
line up several equals signs & \code{align} + \code{\&=} \\
remove one equation number & \code{\textbackslash notag} \\
put ordinary words in math & \code{\textbackslash text\{...\}} \\
make parentheses scale to a fraction & \code{\textbackslash left( ... \textbackslash right)} \\
add a small vertical gap & \code{\textbackslash smallskip} \\
remove one paragraph indent & \code{\textbackslash noindent} \\
show \code{\textbackslash qquad} literally & \code{\textbackslash verb|\textbackslash qquad|} \\
\bottomrule
\end{tabularx}
\end{center}

\section{Reserved Characters - Fast Lookup}

\begin{center}
\begin{tabular}{cccccccccc}
\toprule
Want & \% & \$ & \& & \_ & \# & \{ & \} & \textasciitilde & \textbackslash \\
\midrule
Type & \verb|\%| & \verb|\$| & \verb|\&| & \verb|\_| & \verb|\#| & \verb|\{| & \verb|\}| & \verb|\textasciitilde| & \verb|\textbackslash| \\
\bottomrule
\end{tabular}
\end{center}

\section{Before You Submit: Ten Checks}

\begin{enumerate}
  \item Compile from a clean state and check for real errors.
  \item Search the PDF for \textbf{??} and unresolved citations.
  \item Check that every figure/table is mentioned in the text.
  \item Check that every figure/table has a useful caption.
  \item Check table notes for sample, standard errors, units, and abbreviations where relevant.
  \item Make sure equation symbols are defined near first use.
  \item Inspect wide equations and tables for margin overflow.
  \item Remove TODO notes and draft comments that should not be public.
  \item Verify the bibliography style and all cited references.
  \item Open the final PDF separately and visually inspect several pages, especially tables and figures.
\end{enumerate}

\section{A First-Week Learning Plan}

\begin{enumerate}
  \item \textbf{Day 1:} compile the minimal document; learn text, sections, and paragraphs.
  \item \textbf{Day 2:} learn inline/display math, subscripts, superscripts, fractions, and Greek letters.
  \item \textbf{Day 3:} learn \code{equation}, \code{align}, labels, and references.
  \item \textbf{Day 4:} create one clean table and insert one figure.
  \item \textbf{Day 5:} create a \code{references.bib} file and practice \code{\textbackslash citet}/\code{\textbackslash citep}.
  \item \textbf{Day 6:} split a small paper into files with \code{\textbackslash input}.
  \item \textbf{Day 7:} intentionally create and fix several errors from the debugging section.
\end{enumerate}

\section{Final Advice}

For an Econ PhD student, the goal is not to become a LaTeX programmer. The useful skill is to become fast at turning research content into a clean, reproducible document.

A reliable workflow is:

\begin{enumerate}
  \item start from a working template;
  \item make one change at a time;
  \item compile frequently;
  \item reuse labels, tables, figures, and commands consistently;
  \item let LaTeX handle numbering and layout whenever possible;
  \item debug systematically instead of randomly changing syntax;
  \item keep a short cheat sheet open while writing.
\end{enumerate}

Most day-to-day academic writing uses a surprisingly small subset of LaTeX. For economics research, the recurring tasks are especially predictable: clean prose, notation, cross-references, results tables, figures, citations, appendices, and imported outputs from statistical software. Once the commands in this handbook become familiar, you can learn specialized packages only when a real research need appears.

\end{document}
